% ---------------------------------------------------------------------
%  A total-variation floor and a rank-three classification for the
%  parity-compressed sector of a log-singular Toeplitz form with a
%  prime-power comb.
%
%  Standalone draft.  Self-contained preamble, no external \input:
%  the paper is deliberately independent of any surrounding document
%  set and uses no vocabulary beyond standard numerical analysis.
%
%  DRAFT STATUS (this revision): Sections 1-4 are written out in full,
%  with complete proofs.  Sections 5-9 carry their statements in final
%  form together with proof sketches; the full write-ups, the no-go
%  table of Section 7 and the measurement tables of Section 8 are
%  completed in a later pass.  The typing table of Section 1.3 records
%  the intended status of every numbered statement, and nothing in it is
%  silently upgraded.
% ---------------------------------------------------------------------
\documentclass[11pt,a4paper]{article}
\usepackage[T1]{fontenc}
\usepackage[utf8]{inputenc}
\usepackage[english]{babel}
\usepackage{lmodern}
\usepackage[a4paper,margin=2.6cm]{geometry}
\usepackage{amsmath,amssymb,amsthm}
\usepackage{mathtools}
\usepackage{booktabs,array,tabularx}
\usepackage{xltabular}
\usepackage{enumitem}
\usepackage{microtype}
\usepackage[hidelinks]{hyperref}
\emergencystretch=3em

\theoremstyle{plain}
\newtheorem{theorem}{Theorem}[section]
\newtheorem{proposition}[theorem]{Proposition}
\newtheorem{lemma}[theorem]{Lemma}
\newtheorem{corollary}[theorem]{Corollary}
\theoremstyle{definition}
\newtheorem{definition}[theorem]{Definition}
\newtheorem{convention}[theorem]{Convention}
\newtheorem{remark}[theorem]{Remark}
\newtheorem{measurement}[theorem]{Measurement}
\newtheorem{problem}[theorem]{Problem}

\newcommand{\R}{\mathbb{R}}
\newcommand{\Z}{\mathbb{Z}}
\newcommand{\N}{\mathbb{N}}
\newcommand{\T}{\mathsf{T}}
\newcommand{\Tm}[1]{T_{#1}}
\newcommand{\lmin}{\lambda_{\min}}
\newcommand{\lmax}{\lambda_{\max}}
\newcommand{\cond}{\operatorname{cond}}
\newcommand{\PSD}{\succeq 0}
\newcommand{\PD}{\succ 0}
\newcommand{\Ah}{\widehat{A}}
\newcommand{\ah}{\widehat{a}}
\newcommand{\MM}{\mathcal{M}}
\newcommand{\TV}{\operatorname{TV}}
\newcommand{\Ac}{\operatorname{Ac}}
\newcommand{\Cv}{\operatorname{Cv}}
\DeclareMathOperator{\diag}{diag}
\DeclareMathOperator{\rank}{rank}
\DeclareMathOperator{\tr}{tr}
\DeclareMathOperator{\sgn}{sgn}
\newcommand{\draftnote}[1]{%
  \par\smallskip\noindent\textit{[Draft note: #1]}\par\smallskip}
\setlist{itemsep=2pt,topsep=3pt}
\newcolumntype{Y}{>{\raggedright\arraybackslash}X}

\title{A total-variation floor and a rank-three classification\\
for the parity-compressed sector of a log-singular\\
Toeplitz form with a prime-power comb}
\author{Stefan Hamann}
\date{\today}

\begin{document}
\maketitle

\begin{abstract}
\noindent
This paper is about the near-degeneracy behind a Toeplitz spectral floor: a
total-variation bound, three linear functionals, and a priced table of what the
standard unconditional toolbox reaches. The object is the reflection-odd
compression $A$ of a finite Toeplitz form whose lag sequence is a logarithmic
profile minus an explicit comb, and the question is how small the Rayleigh
quotient of $A$ can be on the low end of the parity spectrum. Four results are
proved. First, a \emph{total-variation floor}: every trial vector whose leading
parity coordinate is normalised to one has lag weights of total variation at
least $1/\mu_1$, where $\mu_1=4\sin^2(\pi/(2h+1))$ is the spectral gap of the
parity compression of the second difference; the floor is closed, involves no
comb node, no taper and no truncation, and it prices the entire trial-vector
route at $O(h^{-2})$ for every vector of bounded value. Second, a \emph{gauge
rigidity} statement: the leading-coordinate normalisation pins only the scale of
the trial direction, both the value and the total variation are homogeneous of
degree two, hence their ratio is a function of the ray and takes the same range
in every reparametrisation of the entry spectrum --- no sector change removes
the floor, and a shifted sector conserves it exactly, floor times transfer
being $1/\mu_1$ identically. Third, an exact \emph{rank-three classification}:
the two-by-two raw Gram block of the low parity modes splits as an archimedean
block minus a comb block, the determinant polarises into three terms, and the
comb--comb term is the polarisation of the determinant on symmetric
two-by-two matrices --- a form of rank three and signature $(1,2)$ --- so the
comb enters the determinant only through \emph{three} linear functionals of its
weights, for every window and every truncation. Fourth, a classification of
what this costs: any majorant of the comb--comb term built from a norm of the
weight vector must dominate the supremum of a rank-three quadratic over the
image of that norm ball, which is where the priced routes lose. The remaining
question --- whether the two rows of the block become collinear at a rate
uniform in the window index --- is stated as an open problem about a finite
bilinear form, with the measured evidence, the placement discriminators and the
priced routes reported as measurements on declared surfaces. The arithmetic
enters the theorems only as the source of the comb positions and weights: every
statement proved here holds for an arbitrary finite comb.
\end{abstract}

\section{Introduction}
\label{sec:intro}

\subsection{The problem, and why the direction is hard}
\label{sec:problem}

Let $c=(c_m)_{m=0}^{M-1}$ be a real lag sequence, let $\Tm{M}(c)$ be its
symmetric Toeplitz matrix, and let $A$ be the compression of $\Tm{M}(c)$ to the
reflection-odd sector of the window (Section~\ref{sec:setting}). The question
this paper is about is a lower bound for the Rayleigh quotient of $A$ on the low
end of a fixed orthonormal mode family --- equivalently, an upper bound for the
extent to which the quadratic form $v\mapsto v^{\T}Av$ can cancel along a smooth
direction.

The classical route from a Toeplitz form to a spectral bound is the symbol. If
$\ell\le\sigma$ pointwise then $\Tm{M}(\ell)\preceq\Tm{M}(\sigma)$ by Parseval,
and a positive pointwise minorant of the symbol transfers to a positive floor
for the form \cite{Szego1939,GrenanderSzego1958,KMS1953,Dirichlet1829,Fejer1916}. For the
class considered here that route is unavailable, and not marginally so: the
symbol is a logarithmic profile minus an explicit cosine comb, the comb dips are
numerous and narrow, and on every window of the declared surface of
Section~\ref{sec:numerics} the infimum of the symbol is negative, while the form
restricted to the parity sector is positive. The smallest Rayleigh quotient is
therefore not produced by a small value of the symbol at some frequency; it is
produced by \emph{cancellation inside the quadratic form}, along a direction that
lives in the low modes. A pointwise minorant is by construction blind to which
vector realises the minimum.

The natural substitute is a trial-vector route: exhibit an explicit direction,
evaluate the form on it, and price the evaluation. This paper is a
classification of that route. Its first two results say that the route has a
closed price which no reparametrisation removes; its second two say that the
object which remains after the price is paid is a quadratic polynomial in three
linear functionals of the comb weights, and that this is exactly why the
unconditional toolbox --- size budgets, symbol arguments, sectors and gauges,
splits, perturbation theory, the bilinear sieve box --- does not reach it. What
is \emph{not} claimed anywhere is a bound on the arithmetic: the three
functionals are finite sums over the comb, all four results hold for an
arbitrary finite comb, and the one remaining question is stated as an open
problem in Section~\ref{sec:open}.

The companion paper \cite{PaperI} treats the same window form from the opposite
direction --- a certified lower bound for a Galerkin defect on a nested ladder
--- and we cite it for the window form and the comb class rather than repeating
them.

\subsection{Contribution}
\label{sec:contrib}

\begin{enumerate}[label=(C\arabic*),leftmargin=3.1em]
\item\label{C:floor} \textbf{The total-variation floor, and the price it
  fixes} (Theorem~\ref{thm:tvfloor} with
  Corollaries~\ref{cor:price}--\ref{cor:abis}). For every trial vector whose
  leading parity coordinate is normalised to one, the lag weights $w$ of the
  associated window vector satisfy
  $\TV(w)\ge|w_0|=\lVert v\rVert^{2}\ge1/\mu_1\ge N^{2}/(4\pi^{2})$, where
  $\mu_1=4\sin^{2}(\pi/N)$, $N=2h+1$, is the spectral gap of the parity
  compression $L_P$ of the second difference \cite{KMS1953}. The proof is four
  steps and uses no property of the comb whatsoever. The load-bearing
  consequence is a \emph{price}: the ratio of the value to the total variation
  is at most $\mu_1$ times the value, so every trial vector of bounded value
  certifies a ratio $O(h^{-2})$, and the boundary-free summation by parts
  \cite{Abel1826} can certify nothing better.
\item\label{C:gauge} \textbf{Gauge rigidity of the normalisation}
  (Theorem~\ref{thm:gauge} with
  Propositions~\ref{prop:shift}--\ref{prop:nullmode} and
  Corollary~\ref{cor:noescape}). The normalisation pins only the scale of the
  trial direction; value and total variation are homogeneous of degree two;
  hence their ratio is a function of the ray, and its range over admissible
  vectors is the \emph{same} for every entry spectrum. A shifted spectrum does
  flatten the floor, but the transfer factor that returns the bound to the
  original scale conserves the product exactly: floor times transfer is
  $1/\mu_1$ at every shift. So (C1) cannot be reparametrised away, and the
  parity compression is not a convenience --- Proposition~\ref{prop:nullmode}
  shows it is what makes the normalisation well posed at all.
\item\label{C:rank} \textbf{The exact bilinear form and the rank-three
  classification} (Section~\ref{sec:rank}: Propositions~\ref{prop:split} and
  \ref{prop:polarisation}, Theorem~\ref{thm:rank3},
  Corollary~\ref{cor:majorant}). The two-by-two raw Gram block of the two lowest
  parity modes splits exactly as $\Ah_2=B-S$ with $B$ archimedean and $S$ a
  weighted sum of comb reads; the determinant polarises as
  $\det\Ah_2=\det B-D(B,S)+\det S$; and $D$ is the polarisation of $\det$ on
  symmetric two-by-two matrices, a form of rank three and signature $(1,2)$ in
  the coordinates $(X_{11},X_{22},X_{12})$. Consequently the comb--comb term is
  $S_{11}S_{22}-S_{12}^{2}$: a quadratic polynomial in \emph{three} linear
  functionals of the comb weights, for every window, every $h$ and every
  truncation. Corollary~\ref{cor:majorant} turns this into the statement about
  majorants that the priced routes of Section~\ref{sec:open} run into.
\item\label{C:nogo} \textbf{A priced table and one open problem}
  (Section~\ref{sec:open}). What remains after (C1)--(C3) is a single question
  about a finite bilinear form: is the near-collinearity of the two rows of
  $\Ah_2$ bounded uniformly in the window index? It is stated as
  Problem~\ref{prob:R1} in three equivalent forms, with the priced routes, the
  measured rates, the density curve and the placement discriminators reported as
  measurements on declared surfaces --- never as an approach to, a weakening of,
  or a consequence of any conjecture.
\end{enumerate}

Two of these four are theorems; two are classifications supported by theorems
together with declared measurements. Which is which is recorded in
Section~\ref{sec:typing} rather than left to the reader.

\subsection{Typing of the results}
\label{sec:typing}

Every numbered statement of this paper carries exactly one of three types, and
the type is part of the statement.

\begin{description}[leftmargin=2.2em,style=nextline]
\item[\textnormal{\emph{proved}}] A written proof is given (or, for classical
  facts, cited). The statement holds for every window and every finite comb
  within its stated hypotheses. Where a hypothesis is itself of a different
  type --- positive definiteness of $A$ on the parity sector, for instance ---
  the statement is \emph{conditional} and says so.
\item[\textnormal{\emph{certified per window}}] The statement is verified on
  each window of a declared finite surface by a computation whose failure mode
  is loud (a completed Cholesky factorisation, an exact identity residual
  against a declared floor). No quantifier over windows is claimed.
\item[\textnormal{\emph{measured}}] The statement is a number, a fitted
  exponent, or a controlled experiment on a declared finite surface, reported as
  such. Fitted exponents are never used as bounds.
\end{description}

{\small
\begin{xltabular}{\textwidth}{@{}llY@{}}
\caption{Typing of every numbered statement. ``Conditional'' marks a proved
implication whose hypothesis is of a weaker type; the hypothesis is then named
in the statement.}
\label{tab:typing}\\
\toprule
statement & type & remark \\
\midrule
\endfirsthead
\toprule
statement & type & remark \\
\midrule
\endhead
\midrule
\multicolumn{3}{r@{}}{\emph{continued on the next page}}\\
\endfoot
\bottomrule
\endlastfoot
\multicolumn{3}{@{}l}{\emph{Section~\ref{sec:setting} --- setting and
  scaffolding}}\\
L1 interlacing (Lem.~\ref{lem:interlacing}) & proved (cited) &
  Cauchy \cite{Cauchy1829}, Courant--Fischer \cite{CourantFischer1920} \\
L2 Schur nesting (Lem.~\ref{lem:nesting}) & proved &
  Schur \cite{Schur1917}, Haynsworth \cite{Haynsworth1968} \\
L3 constrained minimum (Lem.~\ref{lem:constrained}) & proved &
  requires $G_k\PD$ \\
gauge law (Lem.~\ref{lem:gauge}) & proved & $\sum_dw_d=0$, unconditional \\
Abel swap (Lem.~\ref{lem:abel}) & proved & boundary-free, $w_M=0$ \\
C5 cascade forms (Prop.~\ref{prop:cascade}) & proved & prefix / minor /
  regression \\
C6 diagonal invariance (Lem.~\ref{lem:diaginv}) & proved & one line \\
B4 cascade direction (Prop.~\ref{prop:direction}) & \emph{conditional}
  & proved given $\Ah\PD$; the hypothesis is certified per window \\
\midrule
\multicolumn{3}{@{}l}{\emph{Section~\ref{sec:tvfloor} --- the
  total-variation floor}}\\
A the floor (Thm.~\ref{thm:tvfloor}) & proved & comb-free, taper-free \\
A1 the price (Cor.~\ref{cor:price}) & proved & the load-bearing corollary \\
A2 Abel/H\"older (Cor.~\ref{cor:holder}) & proved & \\
A-bis (Cor.~\ref{cor:abis}) & proved & no normalisation needed \\
A4 negative control (Rem.~\ref{rem:negcontrol}) & proved & one line \\
A3 sharpness (Meas.~\ref{meas:sharpness}) & measured & $\mu_1\TV\in[5.00,23.20]$
  on Surface~II \\
\midrule
\multicolumn{3}{@{}l}{\emph{Section~\ref{sec:gauge} --- gauge rigidity}}\\
B homogeneity (Thm.~\ref{thm:gauge}) & proved & \\
B1 floor $\times$ transfer (Prop.~\ref{prop:shift}) & proved & pure shift \\
B2 the null mode (Prop.~\ref{prop:nullmode}) & proved & closed spectra \\
no escape (Cor.~\ref{cor:noescape}) & proved & \\
sector battery (Meas.~\ref{meas:battery}) & certified per window &
  $35\times6\times27$, residual $\le8.8\cdot10^{-10}$ \\
\midrule
\multicolumn{3}{@{}l}{\emph{Section~\ref{sec:twodim} --- the two-dimensional
  reduction}}\\
C1 exact $K=2$ (Prop.~\ref{prop:K2}) & proved & \\
C2 pivot identity (Prop.~\ref{prop:pivot}) & proved & \\
C3 union identity (Prop.~\ref{prop:union}) & proved & \\
C7 Wronskian (Lem.~\ref{lem:wronskian}) & proved & the corner entry $3$ closes
  it \\
C8 maximal minors (Lem.~\ref{lem:minors}) & proved & Lagrange
  \cite{Lagrange1773} \\
C9 determinant reading (Lem.~\ref{lem:detread}) & proved & uses no positivity \\
C10 Gershgorin (Lem.~\ref{lem:gershgorin}) & proved & uniform in the window \\
C11 det-collapse (Prop.~\ref{prop:collapse}) & proved / \emph{conditional} &
  the closed form is unconditional; the harmonic-mean bound needs
  $\Ah_2\PD$ \\
\midrule
\multicolumn{3}{@{}l}{\emph{Section~\ref{sec:rank} --- bilinear form and rank}}\\
D0 exact split (Prop.~\ref{prop:split}) & proved & needs an admissible
  resolution \\
D1 polarisation (Prop.~\ref{prop:polarisation}) & proved & \\
D Theorem D, rank three (Thm.~\ref{thm:rank3}) & proved & every $h$, every
  truncation \\
D3 majorants (Cor.~\ref{cor:majorant}) & proved & declared class of
  majorants \\
D4 Type II rank (Meas.~\ref{meas:typeII}) & measured & bandwidth
  $\Rightarrow\eta$-rank lemma: open writing task \\
D5 reference window (Meas.~\ref{meas:refwindow}) & measured & the
  three-term cancellation \\
\midrule
\multicolumn{3}{@{}l}{\emph{Section~\ref{sec:open} --- the open problem}}\\
Problem~\ref{prob:R1} & open & three equivalent forms \\
rates (Meas.~\ref{meas:rate}) & measured & fitted exponents, declared
  surfaces \\
priced routes (Meas.~\ref{meas:routes}) & proved per route / measured
  exponent & each bound unconditional, each exponent a fit \\
density curve (Meas.~\ref{meas:density}) & measured & undecidability is part
  of the claim \\
placement (Meas.~\ref{meas:placement}) & measured & controlled experiment \\
\end{xltabular}}

Table~\ref{tab:typing} is the first thing to read. Three entries deserve
emphasis. Proposition~\ref{prop:direction} and the harmonic-mean half of
Proposition~\ref{prop:collapse} are \emph{conditional}: they are proved
implications whose hypothesis is positive definiteness of the raw Gram block,
which on the surfaces of Section~\ref{sec:numerics} is a numerical fact about
finitely many finite matrices and nothing more. Measurement~\ref{meas:typeII} is
the one place where a statement of the source material does not convert:
``effective rank $O(1)$'' is not a theorem, the substitute (a bandwidth
$\Rightarrow$ $\eta$-rank lemma with a declared truncation error) is not written
in this revision, and the row therefore stays measured.

\subsection{Notation}
\label{sec:notation}

Notation is fixed once, in Table~\ref{tab:notation}, and not varied. Three
choices need a word. The symbol $\Pi$ denotes the reflection-odd isometry, so
that $B$ can be reserved for the archimedean Gram block of
Section~\ref{sec:rank}; $\MM$ (calligraphic) denotes $\diag(\mu_1,\dots,\mu_K)$,
so that $M$ (upright) keeps its meaning as the window dimension; and the comb
weights, written $\mu_j$ in \cite{PaperI}, are written $\varrho_j$ here, because
$\mu_k$ is needed for the parity eigenvalues. The archimedean kernel is written
$a(\cdot)$ and always carries its argument; the coefficient vector of a trial
vector is written $a$ and never does.

\begin{table}[t]
\centering\small
\begin{tabularx}{\textwidth}{@{}lY@{}}
\toprule
symbol & meaning \\
\midrule
$D,\ M\in2\N,\ \alpha=MD/2,\ h=M/2,\ N=2h+1$ & cell width, window dimension,
  half-width, sector dimension, parity period \\
$\Tm{M}(c)=(c_{|r-r'|})_{r,r'=0}^{M-1}$ & symmetric Toeplitz matrix of the lag
  sequence $c$ \\
$\Pi\in\R^{M\times h}$, $\Pi e_r=\tfrac1{\sqrt2}(e_r-e_{M-1-r})$ &
  reflection-odd isometry, $\Pi^{\T}\Pi=I_h$ \\
$A=\Pi^{\T}\Tm{M}(c)\Pi$, $A_{rr'}=c_{|r-r'|}-c_{M-1-r-r'}$ & the odd window
  form \\
$a(\cdot)$; $S_D(s,u)$; $\varrho_j,u_j$ & archimedean kernel; symmetrised hat;
  comb weights and node positions \\
$L_P$ & odd compression of the lag sequence $(2,-1,0,\dots)$: tridiagonal
  $2I-E-E^{\T}$ with corner entry $3$ \\
$\mu_k=4\sin^{2}(k\pi/N)$, $(t_k)_r=\tfrac2{\sqrt N}\sin\tfrac{2\pi k(r+1)}N$ &
  parity eigenvalues and orthonormal parity modes, $k=1,\dots,h$ \\
$\ah_{ij}=t_i^{\T}At_j$, $\Ah_K=(\ah_{ij})_{i,j\le K}$ & the raw Gram block \\
$\MM=\diag(\mu_1,\dots,\mu_K)$, $G=\MM^{-1/2}\Ah_K\MM^{-1/2}$ & the normalised
  $K$-block \\
$G_k$, $g_k=e_1^{\T}G_k^{-1}e_1$, $s=\mu_1t_1^{\T}A^{-1}t_1$ & leading blocks,
  the ladder, the entry functional \\
$x$ admissible: $x\in\R^{K}$, $x_1=1$ & trial vector \\
$a_k=x_k/\sqrt{\mu_k}$, $v=\sum_{k\le K}a_kt_k$ & coefficient vector and window
  vector \\
$Q(x)=x^{\T}Gx=v^{\T}Av$ & the value \\
$w_0=\Ac_0(v)$, $w_d=2\Ac_d(v)-\Cv_{M-1-d}(v)$ $(d\ge1)$, $w_M:=0$ & lag
  weights; $\Ac$ autocorrelation, $\Cv$ self-convolution \\
$(\Delta w)_d=w_d-w_{d+1}$, $\TV(x)=\lVert\Delta w\rVert_1$ & lag increments and
  total variation \\
$C_d=\sum_{e\le d}c_e$, $C_{-1}:=0$ & partial sums of the lag sequence \\
$B$, $S$, $r_{12}^{2}=\ah_{12}^{2}/(\ah_{11}\ah_{22})$, $\nu_1\ge\nu_2$ &
  archimedean block, comb block, squared correlation, eigenvalues of $\Ah_2$ \\
$D(P,Q)=P_{11}Q_{22}+P_{22}Q_{11}-2P_{12}Q_{12}$ & polarisation of $\det$ on
  symmetric $2\times2$ matrices \\
\bottomrule
\end{tabularx}
\caption{Notation, fixed for the whole paper.}
\label{tab:notation}
\end{table}

$X\PSD$ and $X\PD$ denote positive semidefiniteness and definiteness,
$\preceq$ the L\"owner order, $\lVert\cdot\rVert$ the Euclidean norm,
$e_k$ the $k$-th coordinate vector, $|X|$ the entrywise absolute value of a
matrix, and $\sgn$ the entrywise sign. For $v\in\R^{h}$ the autocorrelation and
self-convolution are
\begin{equation}
\label{eq:accv}
  \Ac_d(v)=\sum_{r}v_rv_{r+d}\quad(d\ge0),
  \qquad
  \Cv_j(v)=\sum_{r+r'=j}v_rv_{r'},
\end{equation}
both sums over indices in $\{0,\dots,h-1\}$ and empty sums being zero.
\emph{Admissible} always means $x_1=1$.

\section{Setting}
\label{sec:setting}

\subsection{The odd window form and the comb class}
\label{sec:window}

We use the window form and the comb class of \cite{PaperI} verbatim and recall
only what is needed. Fix a cell width $D>0$ and an even $M$, put $\alpha=MD/2$,
$h=M/2$, and let $\Pi$ be the reflection-odd isometry of
Table~\ref{tab:notation}. Given a real lag sequence $c=(c_m)_{m=0}^{M-1}$ the
odd window form is
\begin{equation}
\label{eq:Adef}
  A=\Pi^{\T}\Tm{M}(c)\,\Pi\in\R^{h\times h},
  \qquad
  A_{rr'}=c_{|r-r'|}-c_{M-1-r-r'} .
\end{equation}
The class we verify on is the comb class of \cite[Def.~2.6]{PaperI}: with
$a(\cdot)$ the archimedean kernel of that definition, a finite node set
$\{(u_j,\varrho_j)\}_j$ with $u_j>0$, $\varrho_j>0$, and
$S_D(s,u)=\tfrac12[\Delta_D(s-u)+\Delta_D(s+u)]$ the symmetrised hat built from
the unit-height triangle $\Delta_D(t)=\max\{0,1-|t|/D\}$ of half-width $D$,
\begin{equation}
\label{eq:lags}
  c_m=a(mD)-\sum_{j:\,u_j\le2\alpha}\varrho_j\,S_D(mD,u_j),
  \qquad m=0,\dots,M-1 .
\end{equation}
The \emph{prime-power instance} is $u_j=\log n_j$ over the prime powers
$n_j=p^{\nu}$ in increasing order with $\varrho_j=2\Lambda(n_j)/\sqrt{n_j}$,
$\Lambda$ the von Mangoldt weight. By linearity of $c\mapsto A$ the split of
\eqref{eq:lags} induces a split
\begin{equation}
\label{eq:Asplit}
  A=A_{\mathrm{arch}}-A_{\mathrm{comb}},
  \qquad
  A_{\mathrm{comb}}=\sum_{j}\varrho_j\,A^{(j)},
  \quad
  A^{(j)}=\Pi^{\T}\Tm{M}\bigl(S_D(\cdot D,u_j)\bigr)\Pi .
\end{equation}
Throughout, a \emph{resolution} $D$ is called \emph{admissible} if
$D<\min_ju_j$; for the prime-power instance this is implied by the admissibility
convention of \cite[Def.~5.1]{PaperI}, which fixes $D$ as a fraction of the
smallest log-gap in the run.

\begin{remark}[the arithmetic is not used]
\label{rem:combfree}
Theorems~\ref{thm:tvfloor}, \ref{thm:gauge} and \ref{thm:rank3}, and every
lemma, proposition and corollary of Sections~\ref{sec:setting}--\ref{sec:rank},
hold for an \emph{arbitrary} finite comb: the node positions $u_j$ and the
weights $\varrho_j$ enter only as the data of \eqref{eq:lags}, and no property of
the sequence $(u_j)$ or $(\varrho_j)$ is used. Where the prime-power instance is
mentioned it is as the instance the measurements of
Sections~\ref{sec:open}--\ref{sec:numerics} were taken on, and the only
arithmetic facts consumed anywhere are the two classical gap inequalities used
in \cite{PaperI} to check that a resolution is admissible
\cite{Chebyshev1852,RosserSchoenfeld1962}. No arithmetic statement is claimed.
\end{remark}

\subsection{The parity spectrum}
\label{sec:parity}

Let $c^{L}=(2,-1,0,\dots,0)$ and let $L_P$ be the odd window form
\eqref{eq:Adef} of $c^{L}$.

\begin{lemma}[the parity compression and its spectrum]
\label{lem:kms}
$L_P$ is the tridiagonal matrix $2I-E-E^{\T}$ on $\R^{h}$ with corner entry
$(L_P)_{h-1,h-1}=3$; equivalently, $L_P$ is the second-difference form with a
Dirichlet condition at the outer end and the odd reflection $u_h=-u_{h-1}$ at
the inner end. Its eigenpairs are
\begin{equation}
\label{eq:kms}
  \mu_k=4\sin^{2}\Bigl(\frac{k\pi}{N}\Bigr),
  \qquad
  (t_k)_r=\frac2{\sqrt N}\,\sin\frac{2\pi k(r+1)}{N},
  \qquad k=1,\dots,h,\ r=0,\dots,h-1,
\end{equation}
with $N=2h+1$, and $\{t_k\}_{k=1}^{h}$ is an orthonormal basis of $\R^{h}$. In
particular $\mu_1=4\sin^{2}(\pi/N)$ is the smallest eigenvalue and
\begin{equation}
\label{eq:gap}
  \frac1{\mu_1}\;\ge\;\frac{N^{2}}{4\pi^{2}}\;>\;\frac{h^{2}}{\pi^{2}} .
\end{equation}
\end{lemma}

\begin{proof}
For $c^{L}$ the Hankel part of \eqref{eq:Adef} is
$-c^{L}_{M-1-r-r'}$, which is nonzero only if $M-1-r-r'\in\{0,1\}$, i.e.\ only
if $r+r'\in\{2h-2,2h-1\}$; since $r,r'\le h-1$ the only possibility is
$r=r'=h-1$, contributing $-c^{L}_1=+1$. Hence $L_P$ is tridiagonal with
diagonal $2$ except for the corner entry $2+1=3$, which is the first assertion.
Reading the rows as a recurrence, row $0$ is $2u_0-u_1=\mu u_0$, i.e.\ the
convention $u_{-1}=0$, and row $h-1$ is $-u_{h-2}+3u_{h-1}=\mu u_{h-1}$, i.e.\
$-u_{h-2}+2u_{h-1}-u_h=\mu u_{h-1}$ with $u_h=-u_{h-1}$.

For the eigenpairs, put $u_r=\sin(\omega(r+1))$. Then $u_{-1}=0$ and the
interior recurrence holds with $\mu=2-2\cos\omega=4\sin^{2}(\omega/2)$. The
inner condition $u_h=-u_{h-1}$ reads
$\sin(\omega(h+1))+\sin(\omega h)=2\sin\bigl(\tfrac{\omega N}2\bigr)
\cos\bigl(\tfrac\omega2\bigr)=0$, and since $0<\omega<\pi$ the second factor is
positive, so $\omega N/2=k\pi$, i.e.\ $\omega=2k\pi/N$ and
$\mu=4\sin^{2}(k\pi/N)$. The values $k=1,\dots,h$ give $h$ distinct eigenvalues,
hence all of them. For the normalisation,
$\sum_{r=0}^{N-1}\sin^{2}(2\pi kr/N)=N/2$, the term $r=0$ vanishes, and the
remaining $2h$ terms pair up as $r\leftrightarrow N-r$ with equal values, so
$\sum_{r=1}^{h}\sin^{2}(2\pi kr/N)=N/4$ and the factor $2/\sqrt N$ normalises.
Orthogonality is symmetry of $L_P$ together with simplicity of the spectrum.
Finally $\sin t\le t$ on $[0,\pi/2]$ gives
$\mu_1\le4\pi^{2}/N^{2}$, which is \eqref{eq:gap}; the last inequality is
$N=2h+1>2h$.
\end{proof}

The spectrum \eqref{eq:kms} is the classical Kac--Murdock--Szeg\H{o} spectrum of
the second difference in the reflection-odd sector \cite{KMS1953}, and $\mu_1$
is its spectral gap. Everything in Sections~\ref{sec:tvfloor} and
\ref{sec:gauge} is driven by \eqref{eq:gap} and by nothing else.

\subsection{The raw Gram block, the normalised block, and the ladder}
\label{sec:gram}

Fix $K\le h$ (the verified surfaces use $K=16$; nothing below depends on the
value). With $\{t_k\}$ of Lemma~\ref{lem:kms} put
\begin{equation}
\label{eq:gram}
  \ah_{ij}=t_i^{\T}At_j,
  \qquad
  \Ah_K=(\ah_{ij})_{i,j\le K},
  \qquad
  G=\MM^{-1/2}\Ah_K\MM^{-1/2},
  \qquad \MM=\diag(\mu_1,\dots,\mu_K),
\end{equation}
write $G_k$ for the leading $k\times k$ block of $G$, and
\begin{equation}
\label{eq:ladder}
  g_k=e_1^{\T}G_k^{-1}e_1\quad(k=1,\dots,K),
  \qquad
  s=\mu_1\,t_1^{\T}A^{-1}t_1 .
\end{equation}
The following three facts are classical; we record them in the form used later.

\begin{lemma}[interlacing]
\label{lem:interlacing}
Let $X\in\R^{n\times n}$ be symmetric and $P\in\R^{n\times m}$ an isometry.
Then the eigenvalues of $P^{\T}XP$ interlace those of $X$; in particular
$\lmin(X)\le\lmin(P^{\T}XP)$. Consequently a $K$-truncated Rayleigh minimum is
an upper bound for $\lmin$ and never a lower bound.
\end{lemma}

\begin{proof}
Cauchy's interlacing theorem \cite{Cauchy1829} in the Courant--Fischer form
\cite{CourantFischer1920}.
\end{proof}

\begin{lemma}[Schur nesting]
\label{lem:nesting}
Let $X\PD$ be partitioned as
$\bigl[\begin{smallmatrix}\xi&b^{\T}\\ b&Y\end{smallmatrix}\bigr]$ with
$\xi\in\R$. Then $Y\PD$ and
\begin{equation}
\label{eq:schur}
  \frac1{(X^{-1})_{11}}=\xi-b^{\T}Y^{-1}b,
  \qquad
  (X^{-1})_{11}=\frac{\det Y}{\det X} .
\end{equation}
\end{lemma}

\begin{proof}
The first identity is the Schur complement of the $(1,1)$ pivot
\cite{Schur1917,Haynsworth1968}; the second is Jacobi's formula for a cofactor
of the inverse, i.e.\ Cramer's rule.
\end{proof}

\begin{lemma}[constrained minimum]
\label{lem:constrained}
Let $X\PD$ be $k\times k$. Then
\begin{equation}
\label{eq:constrmin}
  \min\{u^{\T}Xu:\ u\in\R^{k},\ u_1=1\}
  =\frac1{e_1^{\T}X^{-1}e_1},
\end{equation}
attained exactly at $u^{\star}=X^{-1}e_1/(e_1^{\T}X^{-1}e_1)$. Equivalently,
\emph{every} admissible $u$ gives $u^{\T}Xu\ge1/(e_1^{\T}X^{-1}e_1)$: an
explicit trial vector is a certified lower bound for the reciprocal, never an
upper one.
\end{lemma}

\begin{proof}
Write $u=u^{\star}+\delta$ with $\delta_1=0$. Since
$Xu^{\star}=e_1/(e_1^{\T}X^{-1}e_1)$, the cross term is
$2\delta^{\T}Xu^{\star}=2\delta_1/(e_1^{\T}X^{-1}e_1)=0$, so
$u^{\T}Xu=(u^{\star})^{\T}Xu^{\star}+\delta^{\T}X\delta$ and
$(u^{\star})^{\T}Xu^{\star}=(u^{\star})^{\T}e_1/(e_1^{\T}X^{-1}e_1)
=1/(e_1^{\T}X^{-1}e_1)$ because $u^{\star}_1=1$. As $X\PD$ the remainder is
positive unless $\delta=0$.
\end{proof}

\begin{proposition}[three forms of the ladder]
\label{prop:cascade}
Assume $G\PD$ and let $G=LL^{\T}$ be its Cholesky factorisation,
$y=L^{-1}e_1$. Then for every $k\le K$:
\begin{align}
\label{eq:prefix}
  \text{(prefix)}\qquad & g_k=\sum_{j\le k}y_j^{2};\\
\label{eq:minor}
  \text{(minor)}\qquad & \frac1{g_k}
    =\frac{\det\Ah_{[1..k]}}{\mu_1\,\det\Ah_{[2..k]}}
    \qquad(\det\Ah_{[2..1]}:=1);\\
\label{eq:regression}
  \text{(regression)}\qquad & \frac{g_k}{g_1}=\frac1{1-R_k^{2}},
    \qquad R_k^{2}=\frac{b^{\T}G_{HH}^{-1}b}{G_{11}},
\end{align}
where in \eqref{eq:regression} the block $G_k$ is partitioned as
$\bigl[\begin{smallmatrix}G_{11}&b^{\T}\\ b&G_{HH}\end{smallmatrix}\bigr]$ with
$H=\{2,\dots,k\}$. In particular $k\mapsto g_k$ is nondecreasing, strictly
increasing whenever $y_k\ne0$, and $R_k^{2}\in[0,1)$ is the squared multiple
correlation of the first coordinate on the remaining $k-1$.
\end{proposition}

\begin{proof}
\emph{Prefix.} For a Cholesky factorisation the leading $k\times k$ block of $G$
is $L_kL_k^{\T}$ with $L_k$ the leading block of $L$, so
$g_k=e_1^{\T}L_k^{-\T}L_k^{-1}e_1=\lVert L_k^{-1}e_1\rVert^{2}$; and since $L$
is lower triangular, forward substitution gives
$L_k^{-1}e_1=(L^{-1}e_1)_{1:k}=y_{1:k}$. Monotonicity is immediate.

\emph{Minor.} By \eqref{eq:schur},
$g_k=\det G_{[2..k]}/\det G_{[1..k]}$. Substituting
$G=\MM^{-1/2}\Ah_K\MM^{-1/2}$ gives
$\det G_{[1..k]}=\det\Ah_{[1..k]}/\prod_{j\le k}\mu_j$ and
$\det G_{[2..k]}=\det\Ah_{[2..k]}/\prod_{2\le j\le k}\mu_j$, and all factors
cancel except $\mu_1$, which yields \eqref{eq:minor}.

\emph{Regression.} By \eqref{eq:schur}, $1/g_k=G_{11}-b^{\T}G_{HH}^{-1}b
=G_{11}(1-R_k^{2})$, while $g_1=1/G_{11}$; divide.
\end{proof}

\begin{lemma}[diagonal invariance of the ladder gain]
\label{lem:diaginv}
For every positive diagonal $\Lambda=\diag(d_1,\dots,d_K)$ the gain
$g_k/g_1=G_{11}\,(G_k^{-1})_{11}$ is unchanged under $G\mapsto\Lambda G\Lambda$.
\end{lemma}

\begin{proof}
$(\Lambda G\Lambda)_{11}=d_1^{2}G_{11}$ and
$\bigl((\Lambda G\Lambda)_k^{-1}\bigr)_{11}=d_1^{-2}(G_k^{-1})_{11}$; the
product is invariant.
\end{proof}

Lemma~\ref{lem:diaginv} localises the phenomenon: the gain does not see the
$\mu^{-1/2}$ normalisation of \eqref{eq:gram} at all, so whatever makes
$g_K/g_1$ large is a property of the raw Gram block $\Ah_K$ alone, and the
parity spectrum enters it only through the choice of the modes.

\begin{proposition}[direction of the ladder]
\label{prop:direction}
Assume $\Ah_h\PD$, i.e.\ $A\PD$ on the parity sector. Then
\begin{equation}
\label{eq:direction}
  s\;\ge\;g_K\;\ge\;g_1=\frac1{G_{11}},
  \qquad\text{equivalently}\qquad
  \frac1s\;\le\;\frac1{g_K}\;\le\;\frac1{g_1},
\end{equation}
with every ladder increment $g_{k+1}-g_k\ge0$. The hypothesis is used only for
the Cholesky factorisation of Proposition~\ref{prop:cascade} and is, on the
surfaces of Section~\ref{sec:numerics}, certified per window.
\end{proposition}

\begin{proof}
Let $T=[t_1,\dots,t_h]$, so $T$ is orthogonal and $\Ah_h=T^{\T}AT$. Then
$A^{-1}=T\Ah_h^{-1}T^{\T}$, hence
$t_1^{\T}A^{-1}t_1=(\Ah_h^{-1})_{11}$; and with
$\MM_h=\diag(\mu_1,\dots,\mu_h)$ and $G^{(h)}=\MM_h^{-1/2}\Ah_h\MM_h^{-1/2}$
one has $\bigl((G^{(h)})^{-1}\bigr)_{11}=\mu_1(\Ah_h^{-1})_{11}$, i.e.\
$s=g_h$. Now apply \eqref{eq:prefix}: $g_k$ is a partial sum of squares, so
$g_1\le g_K\le g_h=s$.
\end{proof}

\subsection{Lag weights, the gauge law, and the boundary-free swap}
\label{sec:lag}

Let $v\in\R^{h}$ and define the \emph{lag weights} of $v$ by
\begin{equation}
\label{eq:weights}
  w_0=\Ac_0(v),
  \qquad
  w_d=2\Ac_d(v)-\Cv_{M-1-d}(v)\quad(1\le d\le M-1),
  \qquad w_M:=0 ,
\end{equation}
with $\Ac,\Cv$ as in \eqref{eq:accv}, and put $(\Delta w)_d=w_d-w_{d+1}$ and
$\TV=\lVert\Delta w\rVert_1=\sum_{d=0}^{M-1}|(\Delta w)_d|$. For an admissible
$x$ we write $w(x)$, $\TV(x)$ for the weights and total variation of the window
vector $v=\sum_{k\le K}a_kt_k$, $a_k=x_k/\sqrt{\mu_k}$.

\begin{lemma}[the weights represent the form; the gauge law]
\label{lem:gauge}
For every $v\in\R^{h}$ and every lag sequence $c$,
\begin{equation}
\label{eq:represent}
  v^{\T}Av=\sum_{d=0}^{M-1}c_d\,w_d,
  \qquad
  w_0=\lVert v\rVert^{2},
  \qquad
  \sum_{d=0}^{M-1}w_d=0 .
\end{equation}
In particular, for an admissible $x$, $Q(x)=x^{\T}Gx=v^{\T}Av=\sum_dc_dw_d(x)$.
\end{lemma}

\begin{proof}
By \eqref{eq:Adef},
$v^{\T}Av=\sum_{r,r'}v_rv_{r'}c_{|r-r'|}-\sum_{r,r'}v_rv_{r'}c_{M-1-r-r'}$. The
first sum collects the lag $d=|r-r'|$ with multiplicity $1$ for $d=0$ and $2$
for $d\ge1$, i.e.\ it equals $c_0\Ac_0(v)+2\sum_{d\ge1}c_d\Ac_d(v)$. The second
sum collects the index $d=M-1-r-r'$, and
since $r,r'\le h-1$ we have $r+r'\le M-2$, so $d\ge1$ throughout and the
coefficient of $c_d$ is $\Cv_{M-1-d}(v)$. This is the first identity of
\eqref{eq:represent} with the weights \eqref{eq:weights}, and it also shows
$w_0=\Ac_0(v)=\lVert v\rVert^{2}$: the reflection term cannot reach $d=0$.

For the gauge law, sum \eqref{eq:weights}. Over all $d\ge0$,
$\Ac_0(v)+2\sum_{d\ge1}\Ac_d(v)=\bigl(\sum_rv_r\bigr)^{2}$, and as $d$ runs
through $1,\dots,M-1$ the index $j=M-1-d$ runs through $0,\dots,M-2=2h-2$, so
$\sum_{d\ge1}\Cv_{M-1-d}(v)=\sum_{j=0}^{2h-2}\Cv_j(v)
=\bigl(\sum_rv_r\bigr)^{2}$. The two contributions cancel.
\end{proof}

The gauge law has a consequence we use once: since $\sum_dw_d=0$, the value
$\sum_dc_dw_d$ is unchanged if a constant is added to every $c_d$. So the
representation \eqref{eq:represent} determines the form only up to an additive
normalisation of $c$, and any bound obtained from it may be optimised over that
normalisation.

\begin{lemma}[boundary-free summation by parts]
\label{lem:abel}
With $C_d=\sum_{e\le d}c_e$, $C_{-1}=0$, and $w_M=0$,
\begin{equation}
\label{eq:abel}
  \sum_{d=0}^{M-1}c_d\,w_d\;=\;\sum_{d=0}^{M-1}C_d\,(\Delta w)_d
\end{equation}
with no boundary term at either end.
\end{lemma}

\begin{proof}
Substituting $c_d=C_d-C_{d-1}$,
\[
  \sum_{d=0}^{M-1}c_dw_d
  =\sum_{d=0}^{M-1}C_dw_d-\sum_{d=0}^{M-1}C_{d-1}w_d
  =\sum_{d=0}^{M-1}C_dw_d-\sum_{d=0}^{M-2}C_dw_{d+1}
  =\sum_{d=0}^{M-1}C_d(w_d-w_{d+1}),
\]
where the second equality shifted the index and used $C_{-1}=0$, and the third
used $w_M=0$ to absorb the last term.
\end{proof}

This is Abel's summation by parts \cite{Abel1826}. Both boundary terms vanish
for structural reasons and not by a choice: $C_{-1}=0$ by definition of a
partial sum, and $w_M=0$ because both terms of \eqref{eq:weights} vanish at
$d=M$.

\section{The total-variation floor}
\label{sec:tvfloor}

\begin{theorem}[the total-variation floor]
\label{thm:tvfloor}
Let $x\in\R^{K}$ be admissible, i.e.\ $x_1=1$, and let $w=w(x)$ be the lag
weights of $v=\sum_{k\le K}a_kt_k$, $a_k=x_k/\sqrt{\mu_k}$. Then
\begin{equation}
\label{eq:tvfloor}
  \TV(x)\;\ge\;|w_0|\;=\;\lVert v\rVert^{2}\;=\;\lVert a\rVert^{2}
  \;\ge\;\frac1{\mu_1}\;=\;\frac1{4\sin^{2}(\pi/N)}
  \;\ge\;\frac{N^{2}}{4\pi^{2}} .
\end{equation}
The bound involves no comb node, no taper and no truncation, and it holds for
every lag sequence $c$ --- indeed the left- and right-hand sides do not depend
on $c$ at all.
\end{theorem}

\begin{proof}
Four steps.

\emph{(i) $w_0=\lVert a\rVert^{2}$.} By Lemma~\ref{lem:gauge},
$w_0=\lVert v\rVert^{2}$; and $\lVert v\rVert^{2}=\lVert a\rVert^{2}$ because
$\{t_k\}$ is orthonormal (Lemma~\ref{lem:kms}).

\emph{(ii) $\TV\ge|w_0|$.} With $w_M=0$ the increments telescope:
$\sum_{d=0}^{M-1}(\Delta w)_d=w_0-w_M=w_0$. Hence
$|w_0|\le\sum_{d=0}^{M-1}|(\Delta w)_d|=\TV$.

\emph{(iii) $\lVert a\rVert^{2}\ge1/\mu_1$.} Admissibility gives
$a_1=x_1/\sqrt{\mu_1}=1/\sqrt{\mu_1}$, so
$\lVert a\rVert^{2}\ge a_1^{2}=1/\mu_1$.

\emph{(iv) $1/\mu_1\ge N^{2}/(4\pi^{2})$.} This is \eqref{eq:gap}, i.e.\
$\sin t\le t$.
\end{proof}

Three features of \eqref{eq:tvfloor} are worth separating. It is a statement
about the \emph{normalisation}, not about the form: step (iii) is the only place
where admissibility is used, and it is where the whole $h^{2}$ comes from. It is
a statement about the \emph{sector}: step (i) uses orthonormality of the parity
modes and step (iv) uses their spectrum, both of which are the classical
Kac--Murdock--Szeg\H{o} data \cite{KMS1953}. And it is \emph{closed}: no
constant is fitted, no window is excluded, and the right-hand side is an
elementary function of $h$.

\begin{corollary}[the price]
\label{cor:price}
For every admissible $x$,
\begin{equation}
\label{eq:price}
  \frac{|Q(x)|}{\TV(x)}\;\le\;\mu_1\,|Q(x)| .
\end{equation}
Consequently, if a family of admissible trial vectors has values bounded by
$|Q(x)|\le U$ with $U$ independent of $h$, the ratio it certifies is at most
$U\mu_1=O(h^{-2})$.
\end{corollary}

\begin{proof}
By Theorem~\ref{thm:tvfloor}, $\TV(x)\ge1/\mu_1>0$, so
$1/\TV(x)\le\mu_1$; multiply by $|Q(x)|\ge0$. The consequence is
$\mu_1=4\sin^{2}(\pi/N)\le4\pi^{2}/N^{2}$.
\end{proof}

Corollary~\ref{cor:price} is the load-bearing statement of the section. It says
that the pair (value, total variation) cannot be improved by choosing a better
trial vector: the ratio is pinned by the value alone, up to the spectral gap.
Any route whose output is the ratio therefore pays a factor $h^{-2}$ that is a
property of the parity spectrum meeting the normalisation, and not a property of
the comb.

\begin{corollary}[what summation by parts can certify]
\label{cor:holder}
For every admissible $x$,
\begin{equation}
\label{eq:holder}
  |Q(x)|\;\le\;\Bigl(\max_{0\le d\le M-1}|C_d|\Bigr)\,\TV(x),
  \qquad\text{hence}\qquad
  \max_{d}|C_d|\;\ge\;\frac{|Q(x)|}{\TV(x)} .
\end{equation}
Combined with \eqref{eq:price}, a family of bounded value certifies
$\max_d|C_d|\ge$ a quantity that is $O(h^{-2})$, and nothing stronger. By the
gauge law of Lemma~\ref{lem:gauge} the estimate may be applied to $c+\gamma$ for
any $\gamma\in\R$, so $\max_d|C_d|$ may be replaced by
$\inf_{\gamma}\max_d|C_d+\gamma(d+1)|$.
\end{corollary}

\begin{proof}
Lemma~\ref{lem:abel} and H\"older: $|\sum_dC_d(\Delta w)_d|
\le\max_d|C_d|\cdot\lVert\Delta w\rVert_1$. The final sentence is the remark
after Lemma~\ref{lem:gauge}, since replacing $c_d$ by $c_d+\gamma$ leaves
$Q$ unchanged and replaces $C_d$ by $C_d+\gamma(d+1)$.
\end{proof}

\begin{corollary}[the floor at one coordinate]
\label{cor:abis}
For every $v\in\R^{h}\setminus\{0\}$ with $t_1^{\T}v\ne0$,
\begin{equation}
\label{eq:abis}
  \frac{\TV(v)}{(t_1^{\T}v)^{2}}\;\ge\;1,
  \qquad\text{because}\qquad
  \TV(v)\;\ge\;|w_0|=\lVert v\rVert^{2}\;\ge\;(t_1^{\T}v)^{2} .
\end{equation}
No normalisation is required: \eqref{eq:abis} holds for every window vector, and
for admissible $x$ it reduces to \eqref{eq:tvfloor} because
$(t_1^{\T}v)^{2}=a_1^{2}=1/\mu_1$.
\end{corollary}

\begin{proof}
The first two relations are steps (i)--(ii) of Theorem~\ref{thm:tvfloor}, which
used no admissibility; the third is Cauchy--Schwarz with $\lVert t_1\rVert=1$.
\end{proof}

\begin{remark}[negative control: the floor is a property of the normalisation]
\label{rem:negcontrol}
Step (iii) is the only step that uses $x_1=1$, and it is indispensable: for
$\lambda\ne0$ the vector $\lambda x$ has $\TV(\lambda x)=\lambda^{2}\TV(x)$
(both $\Ac$ and $\Cv$ are quadratic in $v$, and $v$ is linear in $x$), so
$\mu_1\TV(\lambda x)=\lambda^{2}\mu_1\TV(x)$ falls below $1$ as soon as
$|\lambda|$ is small enough. An inadmissible vector therefore undercuts the
floor, and recognising it as inadmissible is the whole content of the
observation. This is also the reason Theorem~\ref{thm:gauge} below is a
statement about rays.
\end{remark}

\begin{measurement}[sharpness of the floor]
\label{meas:sharpness}
On Surface~II of Section~\ref{sec:numerics} --- $27$ windows of the frame-A
family, $h=142$ to $1445$, all with $\cond(G)\le10^{12}$ --- the floor of
Theorem~\ref{thm:tvfloor} is tight to about one order of magnitude: the measured
ratio is
\begin{equation}
\label{eq:sharp}
  \mu_1\,\TV(x)\in[5.00,\ 23.20]
\end{equation}
at the optimiser of the declared trial family, while the floor itself takes the
values $2.06\cdot10^{3}$ to $1.77\cdot10^{5}$, growing as $h^{+1.997}$ (a
least-squares fit, reported for orientation and not used anywhere) against the
closed $h^{2}$ of \eqref{eq:gap}. The inadmissible control of
Remark~\ref{rem:negcontrol} returns $\mu_1\TV\approx7\cdot10^{-4}$ at
$\lambda=10^{-2}$, i.e.\ it undercuts the floor by the predicted factor
$\lambda^{2}$.
\draftnote{the declared caps of Surface~II, and the definition of the trial
family whose optimiser \eqref{eq:sharp} is evaluated at, are written out with
Section~\ref{sec:numerics}.}
\end{measurement}

We close the section by saying plainly what the floor does not know about. It
does not know where the comb nodes are, or how many there are, or what their
weights are; it does not know whether the trial vector was tapered, and it does
not know at which $K$ the mode family was truncated. It is a statement about
$\{t_k\}$, about $\mu_1$, and about the normalisation $x_1=1$. The next section
shows that the last of these three cannot be traded away.

\section{Gauge rigidity of the normalisation}
\label{sec:gauge}

Theorem~\ref{thm:tvfloor} raises an obvious question: the floor is $1/\mu_1$
because the normalisation is imposed on the mode with eigenvalue $\mu_1$, so
could a different \emph{entry spectrum} --- a reweighted, shifted or otherwise
reparametrised family --- produce a flatter floor and hence a better ratio? The
answer is no, and it is no for a structural reason.

\begin{theorem}[the normalisation is a gauge]
\label{thm:gauge}
For $a\in\R^{K}$ let $v(a)=\sum_{k\le K}a_kt_k$ and
\begin{equation}
\label{eq:functionals}
  \widetilde Q(a)=v(a)^{\T}Av(a),
  \qquad
  \widetilde{\TV}(a)=\lVert\Delta w(v(a))\rVert_1 .
\end{equation}
Then:
\begin{enumerate}[label=\textup{(\roman*)},leftmargin=2.6em]
\item\label{G:hom} $\widetilde Q$ and $\widetilde{\TV}$ are homogeneous of
  degree $2$: $\widetilde Q(\lambda a)=\lambda^{2}\widetilde Q(a)$ and
  $\widetilde{\TV}(\lambda a)=\lambda^{2}\widetilde{\TV}(a)$ for all
  $\lambda\in\R$.
\item\label{G:ray} Hence the ratio
  $\varrho(a)=\widetilde Q(a)/\widetilde{\TV}(a)$ is constant on rays: it is a
  function of the direction of $a$ alone.
\item\label{G:gauge} Let $\mu'=(\mu'_1,\dots,\mu'_K)$ be \emph{any} positive
  entry spectrum and let $x\mapsto a$, $a_k=x_k/\sqrt{\mu'_k}$, be the
  associated parametrisation. Then the constraint $x_1=1$ describes the affine
  slice $\{a:a_1=1/\sqrt{\mu'_1}\}$, every ray $\{\lambda a:\lambda>0\}$ with
  $a_1>0$ meets that slice exactly once, and consequently
  \begin{equation}
  \label{eq:gaugeequal}
    \bigl\{\varrho(a):a_k=x_k/\sqrt{\mu'_k},\ x_1=1\bigr\}
    \;=\;\bigl\{\varrho(a):a\in\R^{K},\ a_1>0\bigr\}
  \end{equation}
  \emph{independently of $\mu'$}. The entry normalisation is a gauge: the range
  of the ratio is the same in every sector.
\end{enumerate}
\end{theorem}

\begin{proof}
\ref{G:hom} $a\mapsto v(a)$ is linear, and by \eqref{eq:accv} both $\Ac_d$ and
$\Cv_j$ are quadratic forms in $v$, so $w(v(\lambda a))=\lambda^{2}w(v(a))$ by
\eqref{eq:weights}, hence $\Delta w$ and its $\ell^{1}$ norm scale by
$\lambda^{2}$; and $\widetilde Q$ is a quadratic form in $a$.

\ref{G:ray} Immediate from \ref{G:hom}, both numerator and denominator having
the same degree.

\ref{G:gauge} The map $x\mapsto a$ is a linear bijection of $\R^{K}$ taking
$\{x_1=1\}$ onto $\{a_1=1/\sqrt{\mu'_1}\}$. A ray $\{\lambda a:\lambda>0\}$ with
$a_1>0$ meets $\{a_1=\theta\}$ exactly once for every $\theta>0$, namely at
$\lambda=\theta/a_1$. So the slice $\{a_1=1/\sqrt{\mu'_1}\}$ is a
cross-section of the family of all rays with $a_1>0$, whatever the value of
$\mu'_1>0$; by \ref{G:ray} the ratio takes the same set of values on the slice
as on the family of rays, which is \eqref{eq:gaugeequal}. The right-hand side
does not mention $\mu'$.
\end{proof}

Theorem~\ref{thm:gauge} is a fact about degree-two homogeneous functionals, and
we state it modestly for that reason. Its force is entirely in the
consequence: since the accessible range of the ratio is a sector-independent
set, and since Theorem~\ref{thm:tvfloor} bounds that range in the parity
parametrisation, no change of sector can enlarge it.

\begin{corollary}[no sector change removes the $h^{2}$]
\label{cor:noescape}
Let $\mu'$ be any positive entry spectrum and let $x'$ be admissible for it.
Then
\begin{equation}
\label{eq:noescape}
  \frac{|\widetilde Q(a')|}{\widetilde{\TV}(a')}
  \;\le\;\mu_1\;\sup\bigl\{|\widetilde Q(a)|:a\ \text{in the parity slice}\bigr\},
\end{equation}
with $\mu_1=4\sin^{2}(\pi/N)$ the parity gap --- i.e.\ the bound of
Corollary~\ref{cor:price} transfers verbatim, with the \emph{parity} gap, to
every sector. In particular a family of trial vectors with $h$-bounded value in
some sector certifies a ratio $O(h^{-2})$ in that sector as well.
\end{corollary}

\begin{proof}
By \eqref{eq:gaugeequal} the value $\varrho(a')$ is attained by some $a$ in the
parity slice ($\mu'=\mu$), i.e.\ by some admissible $x$; apply
Corollary~\ref{cor:price} to that $x$.
\end{proof}

The next proposition is the quantitative form of Corollary~\ref{cor:noescape}
for the simplest reparametrisation, a rigid shift of the entry spectrum. It is a
two-line computation, and we state it as a proposition rather than dressing it
up: the point is not that the computation is hard but that the product it
computes is \emph{exactly} constant.

\begin{proposition}[floor $\times$ transfer is conserved]
\label{prop:shift}
Let $\tau\ge0$ and let $\mu'_k=\mu_k+\tau$ be the shifted entry spectrum, so
that admissibility in the shifted sector reads $a'_1=1/\sqrt{\mu_1+\tau}$. Then
the shifted floor of Theorem~\ref{thm:tvfloor} is
$\widetilde{\TV}(a')\ge1/(\mu_1+\tau)$, the transfer factor relating the two
normalisations is
\begin{equation}
\label{eq:transfer}
  \theta=\frac{(a'_1)^{2}}{a_1^{2}}=\frac{\mu_1}{\mu_1+\tau},
  \qquad
  \frac1\theta=1+\frac\tau{\mu_1},
\end{equation}
and
\begin{equation}
\label{eq:conserved}
  \underbrace{\frac1{\mu_1+\tau}}_{\text{shifted floor}}\cdot
  \underbrace{\Bigl(1+\frac\tau{\mu_1}\Bigr)}_{\text{transfer}}
  \;=\;\frac1{\mu_1+\tau}\cdot\frac{\mu_1+\tau}{\mu_1}
  \;=\;\frac1{\mu_1}
  \qquad\text{identically, at every shift }\tau .
\end{equation}
A shift does flatten the floor; it does not remove the $h^{2}$, which lives in
the ratio and not in the floor.
\end{proposition}

\begin{proof}
The floor is Theorem~\ref{thm:tvfloor} applied in the shifted parametrisation,
whose step (iii) reads $\lVert a'\rVert^{2}\ge(a'_1)^{2}=1/(\mu_1+\tau)$.
For the transfer, the two admissibility constraints differ only in the value of
the first coefficient, $a_1=1/\sqrt{\mu_1}$ against
$a'_1=1/\sqrt{\mu_1+\tau}$; by Theorem~\ref{thm:gauge}\ref{G:hom} a value
computed at $a'$ is converted to the corresponding value on the same ray in the
unshifted slice by the factor $(a_1/a'_1)^{2}=1/\theta$. The displayed product
is then \eqref{eq:conserved}.
\end{proof}

\begin{proposition}[the null mode, and the other sector]
\label{prop:nullmode}
Let $c^{L}=(2,-1,0,\dots)$ be as in Lemma~\ref{lem:kms}.
\begin{enumerate}[label=\textup{(\roman*)},leftmargin=2.6em]
\item\label{N:sine} The parity family $\mu_k=4\sin^{2}(k\pi/N)$,
  $k=1,\dots,h$, is the restriction to $k\ge1$ of the sine family
  $\{4\sin^{2}(k\pi/N)\}_{k\ge0}$, whose member at $k=0$ is $\mu_0=0$.
\item\label{N:noreps} The parametrisation $a_k=x_k/\sqrt{\mu'_k}$ of
  Theorem~\ref{thm:gauge}\ref{G:gauge} requires $\mu'_k>0$. If a mode family
  containing a null mode is used and the normalisation is imposed on that mode,
  the constraint has \emph{no} representative $a$ of finite norm; if the
  normalisation is imposed elsewhere, step (iii) of
  Theorem~\ref{thm:tvfloor} yields the floor of the normalised mode and the
  null mode contributes nothing. In neither case is a flatter admissible floor
  obtained.
\item\label{N:even} Write $\Pi_+e_r=\tfrac1{\sqrt2}(e_r+e_{M-1-r})$ for the
  reflection-even isometry. The complementary compression
  $\Pi_+^{\T}\Tm{M}(c^{L})\Pi_+$ has spectrum
  $\bigl\{4\sin^{2}\bigl((2k+1)\pi/(2N)\bigr)\bigr\}_{k=0}^{h-1}$, whose
  minimum $4\sin^{2}(\pi/(2N))$ is \emph{strictly smaller} than $\mu_1$. Hence
  no reflection sector of the second-difference form has a lowest eigenvalue
  exceeding $\mu_1$.
\end{enumerate}
\end{proposition}

\begin{proof}
\ref{N:sine} is the formula of Lemma~\ref{lem:kms} read at $k=0$.

\ref{N:noreps} is the definition of the parametrisation together with step (iii)
of Theorem~\ref{thm:tvfloor}, which bounds $\lVert a\rVert^{2}$ below by the
square of the \emph{normalised} coordinate only.

\ref{N:even} For $c^{L}$ the even compression is
$(\Pi_+^{\T}\Tm{M}(c^{L})\Pi_+)_{rr'}=c^{L}_{|r-r'|}+c^{L}_{M-1-r-r'}$, which by
the computation in the proof of Lemma~\ref{lem:kms} is tridiagonal
$2I-E-E^{\T}$ with corner entry $2-1=1$; as a recurrence this is $u_{-1}=0$ and
$u_h=+u_{h-1}$. With $u_r=\sin(\omega(r+1))$ the inner condition reads
$\sin(\omega(h+1))-\sin(\omega h)
=2\cos\bigl(\tfrac{\omega N}2\bigr)\sin\bigl(\tfrac\omega2\bigr)=0$, hence
$\omega N/2=(k+\tfrac12)\pi$ and $\mu=4\sin^{2}\bigl((2k+1)\pi/(2N)\bigr)$ for
$k=0,\dots,h-1$, which are $h$ distinct values, hence all. The minimum is at
$k=0$ and $4\sin^{2}(\pi/(2N))<4\sin^{2}(\pi/N)=\mu_1$ because $\sin$ is
increasing on $[0,\pi/2]$.
\end{proof}

Part~\ref{N:even} is the precise sense in which the parity compression is the
best of the two reflection sectors for this purpose, and part~\ref{N:noreps} is
the sense in which enlarging the mode family removes the statement instead of
improving it. Together with Theorem~\ref{thm:gauge} and
Proposition~\ref{prop:shift} this closes the sector question: reweighting,
shifting, or enlarging the entry spectrum changes neither the accessible range
of the ratio nor the conserved product.

\begin{measurement}[the sector battery]
\label{meas:battery}
Theorem~\ref{thm:gauge} and Proposition~\ref{prop:shift} are proved statements;
the following is their machine verification on a declared battery, reported
because it also checks the implementation of the transport map. Over
$35$ weight families $\times$ $6$ taper widths $\times$ $27$ windows of
Surface~II the ratio $\varrho$ is reproduced across sectors with residual at
most $8.8\cdot10^{-10}$ against a declared round-off horizon of
$2.6\cdot10^{-8}$, and the product \eqref{eq:conserved} is reproduced as an
exact identity at every shift of the declared shift list. Transporting the
$x$-coordinates instead of the ray --- the wrong map --- moves the ratio by a
large factor on the spread-out weight families, which is the mutation control of
the battery.
\end{measurement}

\section{The two-dimensional reduction}
\label{sec:twodim}

\draftnote{Sections~\ref{sec:twodim}--\ref{sec:limits} carry their statements in
final form together with proof sketches. The full write-ups, the priced table of
Section~\ref{sec:open} and the measurement tables of
Section~\ref{sec:numerics} are completed in a later pass; the type of each
statement is recorded in Table~\ref{tab:typing} and nothing is upgraded in the
meantime.}

Sections~\ref{sec:tvfloor} and \ref{sec:gauge} priced the trial-vector route.
This section identifies what is left: at $K=2$ the optimal trial vector is
available in closed form, so the entire question collapses onto a single scalar,
the squared correlation of the two lowest parity modes in the metric of $A$.

\begin{proposition}[the closed vector is exact at $K=2$]
\label{prop:K2}
Assume $G_2\PD$ and put $r_{12}^{2}=G_{12}^{2}/(G_{11}G_{22})
=\ah_{12}^{2}/(\ah_{11}\ah_{22})$. Then $u=(1,-G_{21}/G_{22})$ attains the
constrained minimum of Lemma~\ref{lem:constrained} at $k=2$:
\begin{equation}
\label{eq:K2}
  u^{\T}G_2u=\frac1{g_2}=G_{11}\bigl(1-r_{12}^{2}\bigr),
  \qquad
  G_{11}\,g_2=\frac1{1-r_{12}^{2}}\quad\text{identically}.
\end{equation}
\end{proposition}

\begin{proof}[Proof sketch]
Lemma~\ref{lem:nesting} at $k=2$ gives
$1/g_2=G_{11}-G_{12}^{2}/G_{22}=G_{11}(1-r_{12}^{2})$, and
Lemma~\ref{lem:constrained} identifies the minimiser as
$G_2^{-1}e_1/(G_2^{-1}e_1)_1=(1,-G_{21}/G_{22})$. The second identity is the
first divided by $g_1=1/G_{11}$, i.e.\ the case $k=2$ of
\eqref{eq:regression}.
\end{proof}

\begin{proposition}[pivot identity]
\label{prop:pivot}
Assume $G_k\PD$ and let $x^{\star}$ be the minimiser of
Lemma~\ref{lem:constrained}. Then for every $\delta$ with $\delta_1=0$,
\begin{equation}
\label{eq:pivot}
  (x^{\star}+\delta)^{\T}G_k(x^{\star}+\delta)
  =\frac1{g_k}+\delta^{\T}G_k\delta
  \qquad\text{exactly},
\end{equation}
so the relative excess $\rho=g_k\,\delta^{\T}G_k\delta$ of any explicit
candidate is an exact number and never an estimate.
\end{proposition}

\begin{proof}[Proof sketch]
$G_kx^{\star}=e_1/g_k$ and $\delta_1=0$ annihilate the cross term; this is the
computation already carried out in the proof of Lemma~\ref{lem:constrained}.
\end{proof}

\begin{proposition}[the entrywise perturbation, and one cancellation ratio]
\label{prop:union}
Let $x^{\star}$ be as above, $\sigma=\sgn(x^{\star})$,
$\Theta_k=|x^{\star}|^{\T}|G_k|\,|x^{\star}|$, and for $\varepsilon>0$ let
$E$ be the sign-aligned entrywise perturbation $E_{ij}=\varepsilon|G_{ij}|
\sigma_i\sigma_j$. Then
\begin{equation}
\label{eq:union}
  \bigl|(x^{\star})^{\T}Ex^{\star}\bigr|=\varepsilon\,\Theta_k
  \qquad\text{exactly},
\end{equation}
so the entry threshold at which the perturbation exhausts the value is
$\varepsilon_{\mathrm{ent}}(k)=\rho\,(1/g_k)/\Theta_k$ with $\rho$ of
\eqref{eq:pivot}. At $k=2$ the ratio is closed:
\begin{equation}
\label{eq:union2}
  \Theta_2=G_{11}\bigl(1+3r_{12}^{2}\bigr),
  \qquad
  \frac{1/g_2}{\Theta_2}=\frac{1-r_{12}^{2}}{1+3r_{12}^{2}},
\end{equation}
and $1+3r_{12}^{2}\to4$ as $r_{12}^{2}\to1$.
\end{proposition}

\begin{proof}[Proof sketch]
\eqref{eq:union} is the definition of $E$: the signs align every term, so the
sum of $|x^{\star}_i||x^{\star}_j||G_{ij}|$ is attained without cancellation.
At $k=2$, $|x^{\star}|=(1,|G_{12}|/G_{22})$, whence
$\Theta_2=G_{11}+2G_{12}^{2}/G_{22}+G_{12}^{2}/G_{22}
=G_{11}(1+3r_{12}^{2})$; divide by \eqref{eq:K2}.
\end{proof}

\begin{remark}
\label{rem:onedress}
Proposition~\ref{prop:union} is the reason the three natural readings of the
same obstruction --- a determinant ratio, the scalar $1-r_{12}^{2}$, and an
entrywise perturbation threshold --- are one object rather than three: at $k=2$
they differ by the explicit factor $1+3r_{12}^{2}\in[1,4]$. We state the two
formulae and leave the reading to the reader; ``three dresses of one
inequality'' is an observation about identities, not a theorem.
\end{remark}

\begin{lemma}[Wronskian telescope]
\label{lem:wronskian}
Let $u=t_1$, $v=t_2$ be the two lowest parity modes of Lemma~\ref{lem:kms},
extended by $u_{-1}=v_{-1}=0$ and $u_h=-u_{h-1}$, $v_h=-v_{h-1}$, and put
$W_r=u_{r+1}v_r-u_rv_{r+1}$. Then
\begin{equation}
\label{eq:wronskian}
  (\mu_1-\mu_2)\,u_rv_r=W_{r-1}-W_r\quad(0\le r\le h-1),
  \qquad W_{-1}=W_{h-1}=0,
\end{equation}
and summing \eqref{eq:wronskian} re-derives $t_1^{\T}t_2=0$.
\end{lemma}

\begin{proof}[Proof sketch]
$W_{r-1}-W_r=u_r(v_{r-1}+v_{r+1})-v_r(u_{r-1}+u_{r+1})$, and the recurrences
$u_{r-1}+u_{r+1}=(2-\mu_1)u_r$, $v_{r-1}+v_{r+1}=(2-\mu_2)v_r$ give
\eqref{eq:wronskian}. $W_{-1}=0$ is the Dirichlet end; $W_{h-1}=0$ is the odd
reflection $u_h=-u_{h-1}$, $v_h=-v_{h-1}$, i.e.\ exactly the corner entry $3$ of
$L_P$.
\end{proof}

\begin{lemma}[maximal minors carry no arithmetic]
\label{lem:minors}
With $u=t_1/\sqrt{\mu_1}$, $v=t_2/\sqrt{\mu_2}$ and
$\mathcal W_{rr'}=u_rv_{r'}-u_{r'}v_r$,
\begin{equation}
\label{eq:lagrange}
  \tfrac12\sum_{r,r'}\mathcal W_{rr'}^{2}
  =\lVert u\rVert^{2}\lVert v\rVert^{2}-\langle u,v\rangle^{2}
  =\frac1{\mu_1\mu_2}
  \qquad\text{exactly}.
\end{equation}
\end{lemma}

\begin{proof}[Proof sketch]
Lagrange's identity \cite{Lagrange1773} gives the first equality; the second is
orthonormality of $t_1,t_2$ (Lemma~\ref{lem:kms}).
\end{proof}

\begin{lemma}[the determinant reading, free of positivity]
\label{lem:detread}
With $\nu_1\ge\nu_2$ the eigenvalues of $\Ah_2$,
\begin{equation}
\label{eq:detread}
  1-r_{12}^{2}
  =\frac{\det\Ah_2}{\ah_{11}\ah_{22}}
  =\frac{\nu_1\nu_2}{\ah_{11}\ah_{22}}
  =\frac{\det G_2}{G_{11}G_{22}},
\end{equation}
and no positivity of $A$, of $\Ah_2$ or of $G$ is used.
\end{lemma}

\begin{proof}[Proof sketch]
$\det\Ah_2=\ah_{11}\ah_{22}-\ah_{12}^{2}=\nu_1\nu_2$; divide by
$\ah_{11}\ah_{22}$. The last expression is the same computation after the
diagonal congruence \eqref{eq:gram}, which cancels by
Lemma~\ref{lem:diaginv}.
\end{proof}

Lemma~\ref{lem:detread} is load-bearing in a negative sense, and we say so
plainly: it is what makes the whole reduction free of any positivity input about
$A$. Whether $A$ is positive definite on the parity sector is, on the surfaces
of Section~\ref{sec:numerics}, a numerical fact about finitely many finite
matrices; no statement of this paper is routed through a uniform version of it.

\begin{lemma}[an unconditional bound for the large eigenvalue]
\label{lem:gershgorin}
$\nu_1\le\max(\ah_{11},\ah_{22})+|\ah_{12}|$, closed in the three entries and
uniform in the window.
\end{lemma}

\begin{proof}[Proof sketch]
Gershgorin's theorem \cite{Gershgorin1931} for the symmetric $2\times2$ matrix
$\Ah_2$.
\end{proof}

\begin{proposition}[the determinant returns: the $t^{\star}$ collapse]
\label{prop:collapse}
Assume $\ah_{11},\ah_{22}>0$ and consider the one-parameter family
$u(t)=(t,-1)$. At $t^{\star}=\sqrt{\ah_{22}/\ah_{11}}$ the Rayleigh quotient of
$\Ah_2$ is
\begin{equation}
\label{eq:collapse}
  R(t^{\star})
  =\frac{u(t^{\star})^{\T}\Ah_2u(t^{\star})}{\lVert u(t^{\star})\rVert^{2}}
  =\frac{2\sqrt{\ah_{11}\ah_{22}}\;\det\Ah_2}
        {(\ah_{11}+\ah_{22})\bigl(\sqrt{\ah_{11}\ah_{22}}+\ah_{12}\bigr)}
  \qquad\text{identically},
\end{equation}
so the only member of the family that meets the threshold does so by
reintroducing $\det\Ah_2$. If in addition $\Ah_2\PD$, then
\begin{equation}
\label{eq:harmonic}
  \nu_2\;\le\;\frac{2\det\Ah_2}{\ah_{11}+\ah_{22}}
  \;=\;\frac{2\nu_1\nu_2}{\nu_1+\nu_2}\;\le\;2\nu_2,
\end{equation}
i.e.\ the harmonic mean of the two eigenvalues bounds $\nu_2$ from above and is
tight within a factor $2$.
\end{proposition}

\begin{proof}[Proof sketch]
For \eqref{eq:collapse}: with $\tau=t^{\star}$ the numerator is
$\tau^{2}\ah_{11}-2\tau\ah_{12}+\ah_{22}=2(\ah_{22}-\tau\ah_{12})
=2\tau\bigl(\sqrt{\ah_{11}\ah_{22}}-\ah_{12}\bigr)$ and the denominator is
$\tau^{2}+1=(\ah_{11}+\ah_{22})/\ah_{11}$; multiply numerator and denominator by
$\sqrt{\ah_{11}\ah_{22}}+\ah_{12}$ and use
$\ah_{11}\ah_{22}-\ah_{12}^{2}=\det\Ah_2$. For \eqref{eq:harmonic}, substitute
$\tr\Ah_2=\nu_1+\nu_2$ and $\det\Ah_2=\nu_1\nu_2$; the two inequalities are
$\nu_2\le\nu_1$ and $\nu_1+\nu_2\ge\nu_1$ respectively, both of which use
$\nu_2>0$. Without positivity \eqref{eq:harmonic} is false --- if
$\nu_1>0>\nu_2$ and $\tr>0$ the right-hand side is negative --- which is why
the hypothesis is named.
\end{proof}

\begin{remark}
\label{rem:selfsimilar}
Propositions~\ref{prop:K2}--\ref{prop:collapse} are exact reformulations of one
another, and none of them lowers the accuracy demanded of $\det\Ah_2$: the
determinant is reintroduced by \eqref{eq:collapse}, the scalar $1-r_{12}^{2}$ is
the determinant normalised by \eqref{eq:detread}, and the entry threshold of
\eqref{eq:union2} differs from it by the bounded factor $1+3r_{12}^{2}$. We
record this as an observation about the identities, not as a theorem; the
statement ``every exact reformulation costs the same precision'' is a reading of
a family of identities together with measured accuracy demands, and it belongs
in Section~\ref{sec:open} in that form.
\end{remark}

\section{The exact bilinear form and the rank classification}
\label{sec:rank}

This section writes $\Ah_2$ out in terms of the comb and identifies the exact
algebraic shape of the resulting object. Throughout the section the comb nodes
are indexed by $\ell$, so that $i,j$ stay free for matrix indices.

\begin{proposition}[exact split]
\label{prop:split}
Let $D$ be admissible, $D<\min_\ell u_\ell$, and for a node $u_\ell$ let
$m_\ell\in\N$ and $\theta_\ell\in[0,1)$ be defined by
$u_\ell/D=m_\ell+\theta_\ell$. Put $Y^{(m)}=\Pi^{\T}\Tm{M}(\delta_{m})\Pi$ with
$\delta_m$ the lag sequence supported at the single lag $m$, and $Y^{(m)}:=0$
for $m\ge M$. Then, with $B_{ij}=t_i^{\T}A_{\mathrm{arch}}t_j$ and
\begin{equation}
\label{eq:split}
  X^{(\ell)}_{ij}
  =\tfrac12\Bigl[(1-\theta_\ell)\,t_i^{\T}Y^{(m_\ell)}t_j
  +\theta_\ell\,t_i^{\T}Y^{(m_\ell+1)}t_j\Bigr],
  \qquad
  S=\sum_\ell\varrho_\ell\,X^{(\ell)}
  \qquad(i,j\in\{1,2\}),
\end{equation}
one has $\Ah_2=B-S$ exactly, and each of $S_{11},S_{22},S_{12}$ is a
\emph{linear} functional of the weight vector $(\varrho_\ell)_\ell$.
\end{proposition}

\begin{proof}[Proof sketch]
Linearity of $c\mapsto A\mapsto\Ah$ applied to \eqref{eq:lags} and
\eqref{eq:Asplit}. For the two-point read: $S_D(mD,u_\ell)
=\tfrac12[\Delta_D(mD-u_\ell)+\Delta_D(mD+u_\ell)]$; the reflected hat vanishes
for every $m\ge0$ because $mD+u_\ell\ge u_\ell>D$, and the direct hat is nonzero
only for $m\in\{m_\ell,m_\ell+1\}$, where the unit-height triangle $\Delta_D$
takes the values $1-\theta_\ell$ and $\theta_\ell$ respectively. The factor
$\tfrac12$ is the symmetrisation, and a lag index $\ge M$ falls outside the
window and contributes nothing.
\end{proof}

\begin{proposition}[polarisation of the determinant]
\label{prop:polarisation}
With $D(P,Q)=P_{11}Q_{22}+P_{22}Q_{11}-2P_{12}Q_{12}$,
\begin{equation}
\label{eq:polarisation}
  \det\Ah_2=\det B-D(B,S)+\det S,
\end{equation}
where the three terms are respectively archimedean--archimedean, linear in the
comb weights, and bilinear in them. Moreover
\begin{equation}
\label{eq:doublesum}
  \det S=\sum_{\ell,\ell'}\varrho_\ell\varrho_{\ell'}\,
    \mathcal K(\ell,\ell'),
  \qquad
  \mathcal K(\ell,\ell')=\tfrac12D\bigl(X^{(\ell)},X^{(\ell')}\bigr),
\end{equation}
with $X^{(\ell)}$ the $2\times2$ matrix of \eqref{eq:split}.
\end{proposition}

\begin{proof}[Proof sketch]
Both are instances of Theorem~\ref{thm:rank3}\ref{R:pol} below, applied to
$\Ah_2=B-S$ and to $S=\sum_\ell\varrho_\ell X^{(\ell)}$ respectively; the
second uses
bilinearity of $D$ and $\det S=\tfrac12D(S,S)$
\cite{Lagrange1773,CauchyBinet}.
\end{proof}

\begin{theorem}[rank three]
\label{thm:rank3}
Identify a symmetric $2\times2$ matrix $X$ with its coordinate vector
$y(X)=(X_{11},X_{22},X_{12})\in\R^{3}$ and let
\begin{equation}
\label{eq:J}
  J=\begin{pmatrix}0&1&0\\ 1&0&0\\ 0&0&-2\end{pmatrix}.
\end{equation}
Then:
\begin{enumerate}[label=\textup{(\roman*)},leftmargin=2.6em]
\item\label{R:pol} $\det X=\tfrac12\,y(X)^{\T}Jy(X)$ and
  $D(P,Q)=y(P)^{\T}Jy(Q)$; in particular $D$ is the polarisation of $\det$ and
  $\det(P-Q)=\det P-D(P,Q)+\det Q$.
\item\label{R:rank} $J$ has eigenvalues $\{1,-1,-2\}$: it has \textbf{rank
  three} and \textbf{signature $(1,2)$}.
\item\label{R:kernel} Let $P\in\R^{3\times n}$ have columns $y(X^{(\ell)})$,
  $\ell=1,\dots,n$. Then the kernel of \eqref{eq:doublesum} is
  $\mathcal K=\tfrac12P^{\T}JP$, so
  \begin{equation}
  \label{eq:rank3}
    \rank\mathcal K\;\le\;3
    \qquad\text{for every window, every }h\text{ and every truncation},
  \end{equation}
  and the nonzero spectrum of $P^{\T}JP$ equals the nonzero spectrum of
  $JPP^{\T}$, a $3\times3$ matrix.
\item\label{R:three} Consequently
  $\det S=S_{11}S_{22}-S_{12}^{2}$: the comb--comb term of
  \eqref{eq:polarisation} is a quadratic polynomial in exactly \emph{three}
  linear functionals of the comb weights, and depends on the weight vector
  $(\varrho_\ell)$ through nothing else.
\end{enumerate}
\end{theorem}

\begin{proof}[Proof sketch]
\ref{R:pol} Both sides of the first identity are
$X_{11}X_{22}-X_{12}^{2}$; reading off the bilinear form gives
$y(P)^{\T}Jy(Q)=P_{11}Q_{22}+P_{22}Q_{11}-2P_{12}Q_{12}=D(P,Q)$, and expanding
$\tfrac12(y(P)-y(Q))^{\T}J(y(P)-y(Q))$ gives the third relation.
\ref{R:rank} The upper-left $2\times2$ block $\left(\begin{smallmatrix}0&1\\
1&0\end{smallmatrix}\right)$ has eigenvalues $\pm1$, and the remaining entry is
$-2$.
\ref{R:kernel} $\mathcal K(\ell,\ell')
=\tfrac12y(X^{(\ell)})^{\T}Jy(X^{(\ell')})$ is the
$(\ell,\ell')$ entry of $\tfrac12P^{\T}JP$, whose rank is at most $\rank J=3$;
the
statement about spectra is the standard fact
$\operatorname{spec}(UV)\setminus\{0\}=\operatorname{spec}(VU)\setminus\{0\}$
with $U=P^{\T}$, $V=JP$.
\ref{R:three} is \ref{R:pol} applied to $S$, whose coordinates are the three
linear functionals of Proposition~\ref{prop:split}.
\end{proof}

Part~\ref{R:rank} will strike a reader as elementary --- the determinant of a
$2\times2$ matrix is visibly a quadratic form in its three entries --- and it
is. What the theorem is for is part~\ref{R:three} together with the following
corollary, which is where the elementary observation has teeth.

\begin{corollary}[what a majorant of the comb--comb term can see]
\label{cor:majorant}
Let $\mathcal N$ be any norm on the weight vectors and suppose
$|\det S|\le\Phi(\mathcal N(\varrho))$ for all admissible weight vectors, with
$\Phi$ nondecreasing. Then necessarily
\begin{equation}
\label{eq:majorant}
  \Phi(t)\;\ge\;\sup\bigl\{|y_1y_2-y_3^{2}|:\ y\in\R^{3},\
  y=\bigl(S_{11},S_{22},S_{12}\bigr)\ \text{for some }\varrho\ \text{with }
  \mathcal N(\varrho)\le t\bigr\}.
\end{equation}
In particular, no estimate that reaches $\det S$ through a norm of $\varrho$ can
distinguish two weight vectors with the same image
$(S_{11},S_{22},S_{12})$, and any cancellation between the three terms of
\eqref{eq:polarisation} is invisible to it.
\end{corollary}

\begin{proof}[Proof sketch]
Immediate from Theorem~\ref{thm:rank3}\ref{R:three}: $\det S$ is a function of
the image point only, so a bound valid for all $\varrho$ in a norm ball must
dominate the supremum over the image of that ball.
\end{proof}

\begin{measurement}[Type II blocks: a measured rank collapse]
\label{meas:typeII}
The closed weights of Proposition~\ref{prop:split} see a node only through its
position $u_\ell$, i.e.\ for the prime-power instance only through
$\log n_\ell$; under a factorisation $n=bd$ that is $\log b+\log d$, and a
bounded-frequency function of a sum is a finite-rank kernel. On the declared
surface the Type~II blocks of Vaughan's identity \cite{Vaughan1977} built from
these weights are measured to have $2$ singular values carrying $99.9\%$ of the
mass, against $153$ for a generic kernel of the same size, over the scanned
range of the decomposition parameter; the four-term identity for $\Lambda$ is
verified coefficientwise first.
\draftnote{This row is \emph{measured} and stays measured. ``Effective rank
$O(1)$'' is not a theorem. The provable substitute --- \emph{if} the weight
kernel admits an $\Omega$-band expansion in $\log b+\log d$ with declared
truncation error $\eta$, then the block has $\eta$-rank at most $2\Omega+1$ ---
requires a decay estimate for the closed weight kernel in $\log n$ that is not
written in this revision. Bandwidth $\Rightarrow$ $\eta$-rank lemma: open
writing task.}
\end{measurement}

\begin{measurement}[the reference window: a three-term cancellation]
\label{meas:refwindow}
On the reference window $h=540$, with the comb truncated at the node cap
$2.465\cdot10^{4}$, the three terms of \eqref{eq:polarisation} are
\begin{equation}
\label{eq:refwindow}
  \det B=6.18,
  \qquad
  -D(B,S)=-198.8,
  \qquad
  \det S=192.6,
  \qquad
  \text{sum}=4.22\cdot10^{-3},
\end{equation}
i.e.\ five orders of magnitude of cancellation between three pieces of size
$\approx200$; the target value at that window is $1.10\cdot10^{-5}$.
\draftnote{the window's caps and the definition of the target are written out
with Section~\ref{sec:numerics}.}
\end{measurement}

\section{A measured obstruction, and an open problem}
\label{sec:open}

\begin{problem}[three equivalent forms]
\label{prob:R1}
Is there, for each $\varepsilon>0$, a constant $C_\varepsilon$ such that
\emph{uniformly in the window index},
\begin{equation}
\label{eq:R1}
  1-r_{12}^{2}
  =\frac{\det\Ah_2}{\ah_{11}\ah_{22}}
  \;\le\;C_\varepsilon\,h^{-3+\varepsilon}\,?
\end{equation}
Equivalently, by Lemma~\ref{lem:detread} and Lemma~\ref{lem:minors}:
\begin{enumerate}[label=\textup{(\alph*)},leftmargin=2.6em]
\item\label{R1:rows} \emph{Row collinearity.} Do the two rows of $\Ah_2$ ---
  two explicit finite sums against the closed weights of
  Proposition~\ref{prop:split} --- become collinear at rate $h^{-3+\varepsilon}$,
  the sine of the angle between them being
  $|\det\Ah_2|/(\lVert\text{row}_1\rVert\lVert\text{row}_2\rVert)$?
\item\label{R1:det} \emph{Determinant form.} Does the normalised determinant
  $\det\Ah_2/(\ah_{11}\ah_{22})$ satisfy \eqref{eq:R1}?
\item\label{R1:triple} \emph{Joint near-degeneracy.} Does the triple
  $(S_{11},S_{22},S_{12})$ of Theorem~\ref{thm:rank3}\ref{R:three} approach the
  archimedean block $B$ jointly, at relative precision sharpening like
  $h^{-3}$, in the sense that the three terms of \eqref{eq:polarisation} cancel
  to that order?
\end{enumerate}
The problem is open. It is a question about a finite bilinear form; it is not
formulated as, and must not be read as, an approach to, a weakening of, or a
consequence of any conjecture, and no comparison of strengths is made anywhere
in this paper.
\end{problem}

\begin{measurement}[measured rates]
\label{meas:rate}
On the declared surfaces the left-hand side of \eqref{eq:R1} is measured to fall
as a power of $h$ with fitted exponent between $-2.7$ and $-2.95$ depending on
the surface, and the sine-of-angle reading of Problem~\ref{prob:R1}\ref{R1:rows}
falls at the same rate to within the fit uncertainty. Fitted exponents are fits.
\draftnote{the surfaces, the window counts, the caps, and the exponents with
their jackknife bands are written out in the next pass; nothing here is used as
a bound.}
\end{measurement}

\begin{measurement}[priced routes: a table of what the unconditional toolbox
reaches]
\label{meas:routes}
Each row of the table below is an unconditional bound, evaluated on a declared
surface, together with the exponent it certifies; the bound is proved, the
exponent is measured, and the two are typed separately. The routes are:
Cauchy--Schwarz on $\Ah_2$; the large sieve and duality applied to $\det S$
\cite{MontgomeryVaughan1974}; Montgomery--Vaughan applied to the three linear
functionals of Theorem~\ref{thm:rank3}\ref{R:three}; the exact split
$|\det B|+|D(B,S)|+|\det S|$ of \eqref{eq:polarisation}; and Vaughan's
Type I $+$ Type II decomposition \cite{Vaughan1977}. Alongside them sit the
structural no-gos: the symbol infimum (negative on every window of the declared
surface, so no symbol-positivity argument can deliver a floor); perturbation
theory around the parity bottom mode, which bounds at the scale of the matrix
\cite{Kato1949}; band-local domination; and the $\ell^{1}$ mass of the kernel
$\mathcal K$ against its signed value.
\draftnote{Placeholder for the priced table. Each row needs its route's
hypotheses spelled out and its exponent re-read from a declared surface with
caps; until that is written, no exponent is quoted here. Corollary
\ref{cor:majorant} is the structural reason the norm-based rows lose.}
\end{measurement}

\begin{measurement}[the density curve]
\label{meas:density}
The gap between the rate delivered by the best route of
Measurement~\ref{meas:routes} and the rate $h^{-3+\varepsilon}$ of
Problem~\ref{prob:R1}, measured against the density of comb nodes per lag cell,
falls monotonically over three decades of density; the densest reachable bin is
consistent with zero. A true zero, a
power-law approach and a low plateau are \emph{undecidable} under the finite
sieve that the declared caps permit, and this undecidability is part of the
claim, not a caveat attached to it.
\draftnote{Placeholder for the density curve: the six binned values with their
standard errors, the caps that make the dense corner reachable, the jackknife
correction to the scatter, and the explicit undecidability sentence.}
\end{measurement}

\begin{measurement}[placement discriminators]
\label{meas:placement}
Two controlled experiments locate where the arithmetic actually sits. Replacing
the node positions by sorted uniform samples at the \emph{same} weight multiset
leaves the rank-three collapse of Theorem~\ref{thm:rank3} and every
unconditional bound of Measurement~\ref{meas:routes} numerically almost
unchanged, while the \emph{value} of $1-r_{12}^{2}$ moves by orders of
magnitude; and displacing a few nodes by a small fraction of a lag cell moves
$\log$ of the ratio linearly over several declared decades, with a first-harmonic
prediction that under-predicts by a large factor --- the response is not smooth
in the placement. The theorems of this paper are therefore provably not where
the arithmetic is, and the equidistribution of the relevant cell phases is a
classical Weyl statement \cite{Weyl1916}.
\draftnote{Placeholder: the scramble factors, the interventional slope, the
declared decades and the positive-definiteness threshold.}
\end{measurement}

\begin{remark}[what is not claimed]
\label{rem:notclaimed}
Problem~\ref{prob:R1} is open. Nothing in this paper bounds it, and nothing in
this paper is conditional on it. In particular: no statement here is an
arithmetic statement; the positivity of $A$ on the parity sector is used only
where it is named as a hypothesis, and only as a per-window certificate; and the
comparison of a demanded precision against the strength of any conjecture is
deliberately absent from the whole paper.
\end{remark}

\section{Numerical verification}
\label{sec:numerics}

\draftnote{Section skeleton. The floating-point convention, the two declared
surfaces with their caps, and the certificate inventory are written out in the
next pass; the cross-reference table below is complete as it stands.}

\subsection{Floating-point convention and the two surfaces}
\label{sec:surfaces}

Positivity statements are completed Cholesky factorisations relative to a
declared floor in the sense of Wilkinson and Higham
\cite{Wilkinson1963,Higham2002}: with $u=2^{-53}$ the unit roundoff and
$c_h=(h+1)/(1-(h+1)u)$, a completed factorisation of $X-\delta I$ with
$\delta=c_hu\lVert X\rVert_\infty$ certifies $\lmin(X)>0$, and the cheap
certified bound $\lVert X\rVert_2\le\lVert X\rVert_\infty$ is used for the norm
\cite{Gershgorin1931}. Quantities obtained from an eigenvalue routine are
labelled \emph{measured} and never used as certificates. Diagonally dominant
comparisons, where used, are Ostrowski's \cite{Ostrowski1937}.

Two surfaces are declared and never mixed. \emph{Surface~I} (exact, small): the
deepest nodes admitting a window inside a cap on the lag index, on which every
identity of Sections~\ref{sec:setting}--\ref{sec:rank} is recomputed from
scratch. \emph{Surface~II} (measurement): the frame family with
$h=142$ to $1445$ and $\cond(G)\le10^{12}$, on which
Measurements~\ref{meas:sharpness}, \ref{meas:battery}, \ref{meas:refwindow} and
\ref{meas:rate} are taken.

\subsection{Cross-reference table}
\label{sec:crossref}

{\small
\begin{xltabular}{\textwidth}{@{}llY@{}}
\caption{Every numbered statement and the check that recomputes it. Module
identifiers refer to the verification suite accompanying this work; item
numbers are those of the module documentation.}
\label{tab:crossref}\\
\toprule
statement & check & what the check recomputes \\
\midrule
\endfirsthead
\toprule
statement & check & what the check recomputes \\
\midrule
\endhead
\midrule
\multicolumn{3}{r@{}}{\emph{continued on the next page}}\\
\endfoot
\bottomrule
\endlastfoot
Lem.~\ref{lem:kms} (parity spectrum)
  & \texttt{v555}--\texttt{v559}, parity control
  & $L_Pt_k=\mu_kt_k$ with an orthonormal basis;
    $\Pi^{\T}\Tm{M}(c^{L})\Pi=L_P$ at zero tolerance \\
Lem.~\ref{lem:gauge} (gauge law)
  & \texttt{v555} item~(4); \texttt{v559} link~3
  & $\sum_dw_d=0$; $\sum_dc_dw_d=x^{\T}Gx$ \\
Lem.~\ref{lem:abel} (Abel swap)
  & \texttt{v555} item~(4); \texttt{v559} link~9
  & the boundary-free identity at levels $k=1,2,3$ \\
Prop.~\ref{prop:cascade} (three ladder forms)
  & \texttt{v557} item~(1) I1--I3; \texttt{v559} link~12
  & prefix from one Cholesky; minor form; regression form \\
Lem.~\ref{lem:diaginv} (diagonal invariance)
  & \texttt{v557} item~(2)
  & invariance of $g_K/g_1$ under $G\mapsto\Lambda G\Lambda$ \\
Prop.~\ref{prop:direction} (ladder direction)
  & \texttt{v556} item~(5); \texttt{v559} link~4
  & $s\ge g_K\ge g_1$ with strictly positive increments \\
Thm.~\ref{thm:tvfloor} (TV floor)
  & \texttt{v555} item~(2), steps T1--T3
  & $w_0=\lVert a\rVert^{2}$; $\TV\ge|w_0|$; $\mu_1\TV\ge1$ \\
Cor.~\ref{cor:holder} (Abel/H\"older)
  & \texttt{v555} item~(4); \texttt{v559} link~9
  & $|Q|\le\TV\max_d|C_d|$ \\
Cor.~\ref{cor:abis} (floor at one coordinate)
  & \texttt{v556} item~(4b)
  & $\TV\ge\lVert v\rVert^{2}\ge(t_1^{\T}v)^{2}$ \\
Rem.~\ref{rem:negcontrol} (negative control)
  & \texttt{v555} item~(2), control
  & an inadmissible vector undercuts the floor \\
Thm.~\ref{thm:gauge} (gauge)
  & \texttt{v556} item~(1); \texttt{v559} link~10
  & degree-two homogeneity; ratio equal across sectors \\
Prop.~\ref{prop:shift} (conservation)
  & \texttt{v556} item~(2)
  & floor $\times$ transfer $=1/\mu_1$ at every shift \\
Prop.~\ref{prop:K2} ($K=2$ exact)
  & \texttt{v557} item~(3); \texttt{v559} link~13
  & the closed vector attains $1/g_2$; $G_{11}g_2=1/(1-r_{12}^{2})$ \\
Prop.~\ref{prop:pivot}, \ref{prop:union} (pivot, entrywise)
  & \texttt{v557} item~(4)
  & the pivot identity; $|x^{\star\T}Ex^{\star}|=\varepsilon\Theta_K$ \\
Lem.~\ref{lem:constrained} (trial direction)
  & \texttt{v557} item~(5); \texttt{v559} links~4, 7
  & every admissible trial is a lower bound for $1/g_K$ \\
Lem.~\ref{lem:wronskian}, \ref{lem:minors} (Wronskian, minors)
  & \texttt{v558} item~(1) I3
  & the telescope with $W_{-1}=W_{h-1}=0$;
    $\tfrac12\sum\mathcal W^{2}=1/(\mu_1\mu_2)$ \\
Lem.~\ref{lem:detread} (determinant reading)
  & \texttt{v558} items~(2), (6); \texttt{v559} link~14
  & the three equal expressions, no positivity used \\
Lem.~\ref{lem:gershgorin} (Gershgorin)
  & \texttt{v558} item~(2)
  & $\nu_1\le\max(\ah_{11},\ah_{22})+|\ah_{12}|$ \\
Prop.~\ref{prop:collapse} ($t^{\star}$ collapse)
  & \texttt{v558} item~(3)
  & the closed form \eqref{eq:collapse}; the harmonic-mean bound \\
Prop.~\ref{prop:split}, \ref{prop:polarisation} (split, polarisation)
  & \texttt{v558} item~(4) TH1--TH3; \texttt{v559} link~15
  & $\Ah_2=B-S$; \eqref{eq:polarisation}; the kernel $\mathcal K$ \\
Thm.~\ref{thm:rank3} (rank three)
  & \texttt{v558} item~(5) TH4; \texttt{v559} link~15
  & $J$, its spectrum, $\rank\mathcal K\le3$ \\
Meas.~\ref{meas:typeII} (Type II)
  & \texttt{v558} item~(5) TH5
  & singular-value profile against a generic kernel \\
Cor.~\ref{cor:majorant} (majorants)
  & \texttt{v558} items~(4), (6)
  & the wedge reading; the scramble discriminator \\
Prob.~\ref{prob:R1} (the open problem)
  & \texttt{v559} link~16
  & the \emph{shape} is certified per window; the rate and the
    quantifier stay open \\
Meas.~\ref{meas:routes} (priced routes)
  & \texttt{v559} NG1--NG8
  & each classified route exhibited failing as classified \\
\end{xltabular}}

\subsection{Reproducibility}
\label{sec:repro}

\draftnote{Placeholder: the scripts, their runtimes, the number of checks, and
the statement that every factorised matrix stays under the declared cap.}

\section{Limitations, and what is not claimed}
\label{sec:limits}

\begin{enumerate}[label=\textup{(L\arabic*)},leftmargin=3.2em]
\item\label{L:uniform} \textbf{No uniformity in the window index.} Every
  statement typed \emph{certified per window} in Table~\ref{tab:typing} is
  verified window by window on a declared finite surface, and no quantifier over
  windows is claimed. In particular the hypothesis of
  Proposition~\ref{prop:direction} and the positivity hypothesis of
  \eqref{eq:harmonic} are certificates, not theorems.
\item\label{L:positivity} \textbf{Positivity facts are facts about finitely
  many finite matrices.} That $A$ is positive definite on the parity sector on
  the tested windows is a numerical statement with a declared floor. It is not
  routed through, compared with, or used to suggest anything about any
  criterion; Lemma~\ref{lem:detread} exists precisely so that the reduction
  never needs it.
\item\label{L:fits} \textbf{Fits are fits.} Every exponent in
  Measurements~\ref{meas:sharpness}, \ref{meas:rate} and \ref{meas:density} is a
  least-squares fit on a declared surface, reported for orientation. No fit is
  used as a bound anywhere.
\item\label{L:corner} \textbf{The dense corner is exhausted only up to the
  declared caps.} Measurement~\ref{meas:density} states its own
  undecidability, and that statement is part of the claim.
\item\label{L:D4} \textbf{One statement did not convert.}
  Measurement~\ref{meas:typeII} stays measured: the bandwidth $\Rightarrow$
  $\eta$-rank lemma that would make it a theorem is not written, and until it is,
  ``effective rank $O(1)$'' is a measurement and is labelled one.
\item\label{L:arith} \textbf{The comb's arithmetic is not used by any
  theorem.} Remark~\ref{rem:combfree} is meant literally: replacing the
  prime-power instance by an arbitrary finite comb changes no proof in
  Sections~\ref{sec:setting}--\ref{sec:rank}. What changes is the \emph{value}
  of the three functionals of Theorem~\ref{thm:rank3}\ref{R:three}, and
  Measurement~\ref{meas:placement} is the controlled experiment that says so.
\item\label{L:sections} \textbf{Draft scope.} In this revision
  Sections~\ref{sec:setting}--\ref{sec:gauge} are complete with proofs;
  Sections~\ref{sec:twodim}--\ref{sec:limits} carry final statements with proof
  sketches and placeholders as marked. Table~\ref{tab:typing} is the binding
  record of what is claimed at what strength.
\end{enumerate}

\begin{thebibliography}{99}\small

\bibitem{Abel1826} N.~H.~Abel,
\emph{Untersuchungen \"uber die Reihe
$1+\frac m1x+\frac{m(m-1)}{1\cdot2}x^{2}+\dots$},
J. Reine Angew. Math. \textbf{1} (1826), 311--339.

\bibitem{CauchyBinet} A.-L.~Cauchy and J.~Binet,
\emph{M\'emoire sur les fonctions qui ne peuvent obtenir que deux valeurs
\'egales et de signes contraires},
J. \'Ecole Polytech. \textbf{10} (1815), 29--112.

\bibitem{Cauchy1829} A.-L.~Cauchy,
\emph{Sur l'\'equation \`a l'aide de laquelle on d\'etermine les
in\'egalit\'es s\'eculaires des mouvements des plan\`etes},
Oeuvres compl\`etes, 2\textsuperscript{e} s\'erie, tome 9, 174--195 (1829).

\bibitem{Chebyshev1852} P.~L.~Chebyshev,
\emph{M\'emoire sur les nombres premiers},
J. Math. Pures Appl. \textbf{17} (1852), 366--390.

\bibitem{CourantFischer1920} R.~Courant,
\emph{\"Uber die Eigenwerte bei den Differentialgleichungen der
mathematischen Physik},
Math. Z. \textbf{7} (1920), 1--57.

\bibitem{Dirichlet1829} P.~G.~L.~Dirichlet,
\emph{Sur la convergence des s\'eries trigonom\'etriques},
J. Reine Angew. Math. \textbf{4} (1829), 157--169.

\bibitem{Fejer1916} L.~Fej\'er,
\emph{\"Uber trigonometrische Polynome},
J. Reine Angew. Math. \textbf{146} (1916), 53--82.

\bibitem{Gershgorin1931} S.~Gershgorin,
\emph{\"Uber die Abgrenzung der Eigenwerte einer Matrix},
Izv. Akad. Nauk SSSR Ser. Mat. \textbf{6} (1931), 749--754.

\bibitem{GrenanderSzego1958} U.~Grenander and G.~Szeg\H{o},
\emph{Toeplitz Forms and Their Applications},
University of California Press, Berkeley, 1958.

\bibitem{PaperI} S.~Hamann,
\emph{A certified multilevel lower bound for the Galerkin prediction error of a
log-singular Toeplitz form with a prime-power comb}, companion preprint.

\bibitem{Haynsworth1968} E.~V.~Haynsworth,
\emph{Determination of the inertia of a partitioned Hermitian matrix},
Linear Algebra Appl. \textbf{1} (1968), 73--81.

\bibitem{Higham2002} N.~J.~Higham,
\emph{Accuracy and Stability of Numerical Algorithms}, 2nd ed., SIAM,
Philadelphia, 2002.

\bibitem{KMS1953} M.~Kac, W.~L.~Murdock and G.~Szeg\H{o},
\emph{On the eigenvalues of certain Hermitian forms},
J. Rational Mech. Anal. \textbf{2} (1953), 767--800.

\bibitem{Kato1949} T.~Kato,
\emph{On the convergence of the perturbation method, I},
Progr. Theoret. Phys. \textbf{4} (1949), 514--523.

\bibitem{Lagrange1773} J.-L.~Lagrange,
\emph{Solutions analytiques de quelques probl\`emes sur les pyramides
triangulaires},
Nouv. M\'em. Acad. Roy. Sci. Berlin (1773), 149--176.

\bibitem{MontgomeryVaughan1974} H.~L.~Montgomery and R.~C.~Vaughan,
\emph{Hilbert's inequality},
J. London Math. Soc. (2) \textbf{8} (1974), 73--82.

\bibitem{Ostrowski1937} A.~M.~Ostrowski,
\emph{\"Uber die Determinanten mit \"uberwiegender Hauptdiagonale},
Comment. Math. Helv. \textbf{10} (1937), 69--96.

\bibitem{RosserSchoenfeld1962} J.~B.~Rosser and L.~Schoenfeld,
\emph{Approximate formulas for some functions of prime numbers},
Illinois J. Math. \textbf{6} (1962), 64--94.

\bibitem{Schur1917} I.~Schur,
\emph{\"Uber Potenzreihen, die im Innern des Einheitskreises beschr\"ankt
sind},
J. Reine Angew. Math. \textbf{147} (1917), 205--232.

\bibitem{Szego1939} G.~Szeg\H{o},
\emph{Orthogonal Polynomials},
Amer. Math. Soc. Colloquium Publications \textbf{23}, New York, 1939.

\bibitem{Vaughan1977} R.~C.~Vaughan,
\emph{Sommes trigonom\'etriques sur les nombres premiers},
C. R. Acad. Sci. Paris S\'er. A \textbf{285} (1977), 981--983.

\bibitem{Weyl1916} H.~Weyl,
\emph{\"Uber die Gleichverteilung von Zahlen mod. Eins},
Math. Ann. \textbf{77} (1916), 313--352.

\bibitem{Wilkinson1963} J.~H.~Wilkinson,
\emph{Rounding Errors in Algebraic Processes},
Prentice-Hall, Englewood Cliffs, 1963.

\end{thebibliography}

\end{document}
